% concatenated from Solving_Poission_Equation_Wasserstein_Contraction_arxiv.tex
\documentclass[12pt]{article}

\usepackage{latexsym,amsfonts,amsmath,amssymb,epsfig,tabularx,amsthm,dsfont,mathrsfs}
\usepackage{graphicx}
\usepackage{hyperref}
\usepackage{enumerate}
\usepackage{color}


%allows page breaks in algin
\allowdisplaybreaks


\theoremstyle{plain}

\newtheorem{thm}{Theorem}[section]
\newtheorem{rem}[thm]{Remark}
\newtheorem{alg}{Algorithm}
\newtheorem{lem}[thm]{Lemma}
\newtheorem{prop}[thm]{Proposition}
\newtheorem{cor}[thm]{Corollary}
\newtheorem{example}[thm]{Example}
\newtheorem{ass}[thm]{Assumption}
\newtheorem{defi}[thm]{Definition}


\newcommand{\esup}{\operatornamewithlimits{ess\,sup}}
\newcommand{\einf}{\operatornamewithlimits{ess\,inf}}
\renewcommand{\d}{{\rm d}}
\newcommand{\dint}{{\rm d}}
\newcommand{\abs}[1]{\left\vert #1 \right\vert}
\newcommand{\norm}[1]{\left\Vert #1 \right\Vert}
\newcommand{\supp}[1]{\text{supp}\left( #1 \right)}


\newcommand{\N}{\mathbb{N}}
\newcommand{\E}{\mathbb{E}}
\newcommand{\R}{\mathbb{R}}
\newcommand{\Q}{\mathbb{Q}}
\newcommand{\Z}{\mathbb{Z}}

\newcommand{\cX}{\mathcal{X}}
\newcommand{\cY}{\mathcal{Y}}
\newcommand{\cZ}{\mathcal{Z}}
\newcommand{\cD}{\mathcal{D}}
\newcommand{\keff}{k_{\text{eff}}}

\DeclareMathOperator{\Var}{\text{Var}}
\DeclareMathOperator{\Cov}{\text{Cov}}
\DeclareMathOperator{\eps}{\varepsilon}
\DeclareMathOperator{\med}{\text{med}}
\DeclareMathOperator{\Id}{\text{Id}}
\DeclareMathOperator{\KL}{\text{KL}}
\DeclareMathOperator*{\argmin}{argmin}



\newcounter{Ax}

\title{
Solving Poisson's equation for Wasserstein contractive Markov chains
}

\author{Julian Hofstadler\thanks{Email: jh4272@bath.ac.uk}\\[1ex]University of Bath} 



\date{\today}



\begin{document}
	\maketitle

\begin{abstract}
\noindent
We study Poisson's equation in the context of general state space Markov chains. For chains satisfying a contraction assumption w.r.t.~a Wasserstein distance, we show that a solution exists for Lipschitz functions and investigate its regularity properties. If the kernel is additionally reversible we are also able to show that solutions for $L^p$ functions exist. Combining our findings with Doob's inequalities for martingales we derive maximal inequalities for contractive Markov chains. A number of examples is provided to demonstrate the applicability of our results, in particular in the context of Markov chain Monte Carlo methods. 
\end{abstract}

\noindent
\textbf{MSC classification:} 60J05, 60J22, 65C05

\noindent
\textbf{Keywords:} Wasserstein contraction, Poisson's equation, maximal inequality

\section{Introduction}
Let $(X_n)_{n \in \N}$ be a Markov chain\footnote{Here and later we assume that $(X_n)_{n \in \N}$ is defined on a sufficiently regular underlying probability space $(\Omega, \mathcal{F}, \mathbb{P})$.} evolving in some (measurable\footnote{To be precise, $\cX$ is assumed to be Polish and equipped with the corresponding Borel $\sigma$-algebra, cf. Section \ref{sec:Wasserstein_contractive_chains}.}) state space $\cX$ and having invariant distribution $\pi$. Assume $f \colon \cX \to \R$ is a $\pi$-integrable (hence measurable) function. 
Describing the convergence properties of 
\[
S_n f = \sum_{j=1}^{n}f(X_j),
\]
is a classical task in probability theory. 
Already Kolmogoroff devoted a chapter to this question (for iid sequences) when laying the foundations of modern probability in \cite{kolmogoroff1933grundbegriffe}.
By now the case of iid sequences is well-understood and we know that for any $\pi$-integrable function $f$ we have $\frac{1}{n} S_n f \to \pi(f)$ almost surely, see e.g. \cite{etemadi1981elementary} for an elementary proof. 
Other classical results include the central limit theorems (CLT), see for instance \cite{dudley2002real}, or concentration inequalities, cf. \cite[Appendix B]{pollard1984stochastic}. 

It is possible to generalise these results, under suitable assumptions, to more general Markov chains $(X_n)_{n \in \N}$.
Indeed, (under mild assumptions) we obtain the strong law of large numbers \cite{asmussen2011ergodic}. 
Provided $(X_n)_{n \in \N}$ is sufficiently regular, also CLTs \cite{maigret1978theoreme,douc2018markov}, the law of the iterated logarithm \cite{meyn2009markov,bolt2012invariance} and concentration inequalities \cite{GLYNN2002hoeffding,MIASOJEDOW2014hoeffding} are known for Markov chains.
See also \cite{adamczak2008tail,adamczak2015exponential} for uniform concentration results.


A powerful tool to prove these results, which was indeed used in the aforementioned works \cite{douc2018markov,GLYNN2002hoeffding,meyn2009markov,bolt2012invariance,maigret1978theoreme}, is to find $u \colon \cX \to \R$ which for given $f\colon \cX \to \R$ satisfies the identity  
\begin{equation}\label{equ:Poission_equation}
u(x) - Pu(x) = f(x) - \pi(f).
\end{equation}
Equation \eqref{equ:Poission_equation} is called \textit{Poission's equation} (with \textit{forcing function\footnote{Sometimes $f$ is also referred to as \textit{charge} \cite{nummelin1991poisson} or \textit{second member} \cite{revuz1984markov}.}} $f$), and $u$ is called a \textit{solution} of Poission's equation.
For any such $u$ we have 
\begin{equation}\label{equ:martingale_decomposition}
S_n f = \sum_{j=1}^n\left(  u(X_j) - Pu(X_j)\right) = M_n + u(X_1) -u(X_{n+1}), 
\end{equation}
where $M_n = \sum_{j=1}^n\left(  u(X_{j+1}) - Pu(X_j)\right)$.
If $u(X_{j+1}) - Pu(X_j)$ is an integrable random variable for any $j\in \N$, then, setting $\mathcal{F}_j = \sigma(X_k \colon 1 \leq k \leq j+1)$, we obtain that $(M_n)_{n \in \N}$ is a martingale w.r.t.~$(\mathcal{F}_n)_{n \in\N}$.  

Poisson's equation and a generalised version of \eqref{equ:martingale_decomposition} are also of interest to study the convergence of \textit{adaptive} Markov chain Monte Carlo algorithms, see e.g. \cite{andrieu2006ergodicity,saksman2010ergodicity,atchade2013Bayesian,hofstadler2024convergenceRates,laitinen2024invitation}.
These methods often outperform their non-adaptive counterparts in numerical simulations, which renders them particularly relevant in applications. 
The adaptive behaviour leads to more complex martingale and remainder terms, since one needs to solve Poisson's equation for a whole class of kernels $(P_\gamma)_{\gamma \in \mathcal{I}}$ (for one fixed forcing function $f$ and a measurable parameter space $\mathcal{I}$). 
Thus, in this context it is crucial to have knowledge about the stability, e.g. integrability properties, of the corresponding $u_\gamma$'s.  


Connecting Poisson's equation (in the context of general state space Markov chains) and potential theory goes back to \cite{neveu1972potentiel} and \cite{revuz1984markov}, see also \cite{nummelin1991poisson}.
Existence and regularity of a solution $u$ under appropriate Lyapunov function and drift conditions was studied in \cite{glynn1996liapounov}. 
Recently, the authors of \cite{glynn2025computableboundssolutionpoissons} studied Poisson's equation and provided computable bounds for $u$ under drift and minorisation conditions, cf.~Assumption~2.1 there.
In particular, \cite[Equation (2.10) and Theorem 2.1]{glynn2025computableboundssolutionpoissons} establish a probabilistic representation of $u$.
Similar bounds and connections to perturbation theory of Markov chains/kernels were also shown in \cite{herve2025computable}, where a different proof technique was used (compared to \cite{glynn2025computableboundssolutionpoissons}). 
In addition to that, coupled Markov chains may be used to derive solutions of Poisson's equation \cite{douc2025solvingpoissonequationusing}.
The main idea in this approach is to consider two chains which (almost surely) meet after a finite amount of Markov steps, and allow to compute $u$. 
The authors of \cite{qu2025deeplearningpoisson} investigated connections to AI and developed a method to compute/approximate $u$ via ReLU networks.
The reader may also consult \cite[Section 17.4.1]{meyn2009markov} as well as \cite[Section 21.2]{douc2018markov} for an overview about Poisson's equation. 
Finally, it is worth mentioning that there are results specifically for countable $\cX$, see e.g. \cite{glynn2024solution} and the references therein. 

The goal of this article is to study Poisson's equation for chains satisfying a Wasserstein contraction assumption. 
Additionally, we derive maximal inequalities $S_n f$. 
Our assumptions contain uniformly and $V$-uniformly ergodic chains as special case, see Lemma \ref{lem:uniformly_ergodic_implies_contraction} and Lemma \ref{lem:V_uniformly_ergodic_implies_contraction}, which are e.g.~of interest in the Markov chain Monte Carlo literature.
Examples of chains fitting into this setting can be found in Section \ref{sec:Solutions}. 

Poisson's equation was studied under Wasserstein contraction assumptions in \cite{bolt2012invariance} and recently also in \cite{hofstadler2024convergenceRates}. 
The special cases of uniformly, or $V$-uniformly ergodic chains were studied in \cite{GLYNN2002hoeffding} and \cite{glynn1996liapounov}, respectively. 
Often, e.g. in \cite{nummelin1991poisson,GLYNN2002hoeffding,meyn2009markov,bolt2012invariance,grama2018limit,herve2025computable, hofstadler2024convergenceRates}, computing $u$ for a given forcing function $f$ is done by considering the series
\[
\widetilde{u}(x) = 
\sum_{\ell=0}^\infty\left(  P^\ell f(x ) - \pi(f)\right). 
\]
Under suitable assumptions on the Markov chain $\widetilde{u}$ is a well-defined solution of Poisson's equation, we refer to the aforementioned papers. 
Alternative approaches work with hitting times, cf. \cite[Corollary 2.5]{nummelin1991poisson}, see also the recent work \cite{glynn2025computableboundssolutionpoissons}; or coupled Markov chains \cite{douc2025solvingpoissonequationusing}. 


In the present paper, however, we use spectral properties of Markov operators on suitable Banach spaces containing centred Lipschitz functions. 
Then, one `recovers' $\widetilde{u}$ from a Neumann series argument. 
In Theorem~\ref{thm:Solution_Poission_Lipschitz_general} and Theorem \ref{thm:Solution_Poisson_with_L_p_bound} we prove existence of $u$ and deduce regularity properties, such as $L^p$ bounds. 
To the best of our knowledge a thorough analysis based on this approach was not done so far in the Markov chain Monte Carlo literature. 
Additionally, our results can be transferred to the $L^p$ under a reversibility assumption thereby extending known results, see the discussion above Theorem~\ref{thm:poisson_equation_spectral_gap_general}. 
Using decomposition \eqref{equ:martingale_decomposition} we apply our results to obtain maximal inequalities in Section~\ref{sec:applications}. 
To the best of our knowledge such inequalities are novel in the Markov chain Monte Carlo literature, and complement the results of \cite{ollivier2009ricci, Joulin2010Curvature,paulin2016mixing}.  

Maximal inequalities are bounds for $\max_{1 \leq k \leq n}\vert S_n f - n \pi(f)\vert $.
These provide a powerful tool to study the path-wise behaviour of $S_n f$. E.g. one can derive (quantitative) laws of large numbers by using results such as Theorem~\ref{thm:Doob_maximal_bounded_diameter} or Theorem~\ref{thm:Doob_maximal_inequality_finite_moment_assumption}, and considering (properly renormalised) subsequences $(S_{n_k} f)_{k \in \N}$ rather than $(S_n f)_{n \in \N}$, see Theorem~\ref{thm:almost_sure_convergence_NUTS}.
Furthermore, bounds as in Theorem~\ref{thm:L_2_maximal_inequality_bounded_diameter} or Corollary~\ref{cor:L_2_maximal_inequality_finite_moment_assumption}, can be used to obtain information about the maximal error of the whole simulation, instead of the last step.  
This may be particularly relevant in the context of Markov chain Monte Carlo algorithms. 

The remainder of the paper is organised as follows. 
In Section \ref{sec:Wasserstein_contractive_chains} we provide definitions and fundamental properties of Wasserstein distances, as well as Wasserstein contractive chains. Additionally, a number of examples can be found there. 
Section \ref{sec:Solutions} contains our main results about the existence and regularity concerning solutions of Poisson's equation for Wasserstein contractive chains. 
In Section \ref{sec:applications} we derive maximal inequalities for contractive chains, which are then applied to deduce path-wise convergence properties for the No-U-Turn-Sampler \cite{hoffman2014no}. 


\section{Wasserstein contractive Markov chains}\label{sec:Wasserstein_contractive_chains}
In this section we review concepts from the Markov chain literature and related fields which are required later.
In particular, we provide a definition of \textit{Wasserstein contractive Markov chains}, cf. Definition \ref{def:contractive_chain}, which are our main object to study Poisson's equation. 
At the end of this section we include a number of examples, which fit into our framework.



Throughout all sections $(\cX , \rho)$ is a (non-trivial) Polish space and $\mathcal{B}(\cX)$ is the corresponding Borel $\sigma$-algebra. Moreover, $d \colon \cX \times \cX \to \R^+$ is a generic lower semi-continuous (w.r.t.~the product topology of $\rho$) metric.   



\subsection{Definitions and notation}
Let $(X_n)_{n \in \N}$ be a Markov chain evolving in $\cX$ with transition kernel $P\colon \cX \times \mathcal{B}(\cX) \to [0,1]$. 
Unless specified otherwise, $\nu$ denotes the initial distribution, that is, $\nu$ is a probability measure on $(\cX, \mathcal{B}(\cX))$ and we have $X_1 \sim \nu$.
We follow the usual convention and write $(\Omega, \mathcal{F},\mathbb{P}_\nu)$ and $\E_\nu$ for the underlying probability space and expectation, respectively, cf. \cite{douc2018markov} for details.

Given a probability measure $\mu$ on $(\cX, \mathcal{B}(\cX))$, we define $\mu P$ as 
\[
\mu P (\cdot ) = \int_\cX P(x, \cdot) \mu(\d x),
\] 
which is again a probability measure. Hence $\mu \mapsto \mu P$ is a well-defined mapping from the set of probability measures on $(\cX , \mathcal{B}(\cX))$, write $\mathcal{P(\cX)}$, into itself. 
If $\mu$ and $P$ satisfy 
\[
\int_A P(x, B) \mu(\d x) = \int_B P(x, A) \mu(\d x),
\]
for any measurable sets $A, B \in \mathcal{B}(\cX)$, then $P$ is called \textit{reversible} w.r.t.~$\mu$. In this case $\mu$ is an \textit{invariant distribution}, i.e. $\mu P = \mu$, see \cite{douc2018markov}.  

For a sufficiently regular function $f\colon \cX \to \R$ we define 
\[
Pf (x) = \int_\cX f(x') P(x , \d x').
\]
Typically, $P$ is called \textit{Markov operator} when we work with $Pf$. For more details about $P$ acting on either probability measures or functions, the reader may consult e.g. the books \cite{meyn2009markov, douc2018markov}. 


Let $\mu_1, \mu_2 \in \mathcal{P(\cX)}$ and denote by $\mathcal{C}(\mu_1, \mu_2)$ the set of couplings between $\mu_1$ and $\mu_2$. The 1-Wasserstein distance, sometimes also called Kantorovich-Rubinstein metric, is defined as 
\[
\mathcal{W}(\mu_1, \mu_2) = \inf_{\eta \in \mathcal{C}(\mu_1, \mu_2)} \int_{\cX \times \cX} d(x,y) \eta(\d x , \d y).
\]
If $\mathcal{W}$ is finite, then it indeed satisfies all axioms of a metric, and hence allows to measure the distance between $\mu_1$ and $\mu_2$. It is possible to generalise the above definition to so-called $p$-Wasserstein distances, with $p \geq 1$, which are fundamental objects in optimal transport, we refer to \cite{villani2009optimal}.
However, we only consider the 1-Wasserstein distance in this article, whence we shall simply call it \textit{Wasserstein distance}.

Define the \textit{Kantorovich norm} of a Markov kernel $P \colon \cX \times \mathcal{B}(\cX) \to [0,1]$ as 
\[
\tau(P) = \sup_{x \neq y} \frac{\mathcal{W}(P(x, \cdot ), P(y,\cdot ))}{d(x,y)}.
\]
This quantity was used e.g.~by Dobrushin \cite{dobrushin1956central}, see Equation 1.5 there, who later called it Kantorovich norm in \cite{dobrushin96}.
In \cite{douc2018markov} the terminology \textit{Dobrushin coefficient} is used and in \cite{rudolf2018pertubation} it was called \textit{generalised ergodicity coefficient}.
From \cite{ollivier2009ricci}, see also \cite{Joulin2010Curvature,paulin2016mixing}, it is known that $\tau(P)$ is connected to the Ricci curvature of the corresponding Markov chain.


\begin{defi}\label{def:contractive_chain}
A Markov chain $(X_n)_{n \in \N}$ with transition kernel $P$ is called \textit{Wasserstein contractive}, if $\tau(P)<\infty$ and $\tau(P^m)<1$ for some $m \in \N$. 
\end{defi}

\begin{rem}
From this point onwards, we always assume that $m \in \N$ in Definition \ref{def:contractive_chain} is the minimal $m \in \N$ such that $\tau(P^m)<1$.
\end{rem}

\begin{rem}
The assumption $\tau(P^m)<1$ is closely related to the so-called multistep coarse curvature of \cite{paulin2016mixing}, and whence can be seen as generalisation of the assumptions in \cite{ollivier2009ricci,Joulin2010Curvature}.
\end{rem}
Let us briefly explain the terminology. Propositions 14.3 and 14.4 of \cite{dobrushin96}, see also \cite[Proposition 2.1]{rudolf2018pertubation}; \cite[Proposition 20.3.3]{douc2018markov}, yield 
\begin{itemize}
	\item[(i)] $\tau(P P') \leq \tau (P) \tau (P')$, and 
	\item[(ii)]$\mathcal{W}(\mu_1 P, \mu_2 P) \leq \tau(P) \mathcal{W}(\mu_1, \mu_2)$,
\end{itemize}
where $P, P'$ are Markov kernels and $\mu_1, \mu_2 \in \mathcal{P(\cX)}$.
By (ii), if $\tau(P^m)<1$, then $\mu \mapsto \mu P^m$ (considered as mapping from $\mathcal{P(\cX)}$ to $\mathcal{P(\cX)}$) is a contraction w.r.t.~the Wasserstein distance $\mathcal{W}$. 
Therefore, one might expect that there is a fixed point (aka invariant distribution), say $\pi$. 
Yet, $\mathcal{W}$ might be infinite in general. 
However, if $\int_\cX d(x_0, x) P(x_0, \d x) < \infty$ for some $x_0 \in \cX$, then we can prove the existence of $\pi$, c.f. \cite[Theorem 20.3.4]{douc2018markov}, see also \cite{paulin2016mixing}.

\begin{rem}\label{rem:def_Lambda}
From property (i) above it follows that for a contractive chain
\[
\sum_{n=0}^\infty \tau(P^n) = \sum_{\ell=0}^\infty\sum_{j=0}^{m-1}\tau(P^{m\ell +j}) \leq \frac{m \tau(P)^m}{1-\tau(P^m)} < \infty.
\]
Hence, for any contractive chain we may define $\Lambda := \sum_{n=0}^\infty \tau(P^n) \in [0, \infty)$.
\end{rem}

\subsection{Examples}
In this section we present examples of Wasserstein contractive Markov chains. 
In particular, we present a number of contractive chains that correspond to Markov chain Monte Carlo (MCMC) algorithms, which demonstrates the connection of our setting to applications. 


\begin{example}[{\cite[Example 20.3.5]{douc2018markov}}]
	Let $\cX = [0,1]$, $ d(x,y) = \vert x-y\vert$ and $\mathcal{B}(\cX)$ be the corresponding Borel sets.
	We consider a Markov chain $(X_n)_{n \in \N}$ given by 
	\[
	X_{n+1} = \frac{1}{2}\left( X_n + Z_{n+1} \right),
	\]
	where $(Z_n)_{n \in \N}$ is an iid sequence of Bernoulli random variables with mean $1/2$, and $X_1$ is another random variable on $[0,1]$, independent of $(Z_n)_{n \in \N}$. 
	It is shown in \cite{douc2018markov} that in this case we have $\tau(P) = 1/2$.
\end{example}


\begin{example}[{Heat bath for the Ising model -- cf. \cite[Example 17]{ollivier2009ricci} and \cite[Section 2.2]{Joulin2010Curvature}}]
Here we set $\cX = \{-1,1\}^G$, for some finite graph $G$. We define $\pi$ to be the probability measure which has density 
\[
p(x) \propto \exp \left( - \beta\sum_{g_1 \sim g_2} x(g_1)x(g_2) - h \sum_{g \in G} x(g) \right) = \exp (-U(x)),
\]
where $\beta\geq 0$, $h \in \R$ and $g_1 \sim g_2$ means that $g_1$ and $g_2$ are connected. 
The normalisation constant is given by $Z = \sum_{x \in \cX} \exp(-U(x))$. 

Set $d(x_1,x_2) = \frac{1}{2}\sum_{g \in G} \vert x_1(g)- x_2(g) \vert$. Then, $d(x_1, x_2)$ measures to the number of $g \in G$ for which $x_1(g)$ and $x_2(g)$ are different.  
Moreover, for $x \in \cX$ and $g \in G$ we define $y_-\in \cX$ and $y_+ \in \cX$ by setting $y_-(g)=-1$ and $y_+(g)=1$, and $y_-(g') = y_+(g') = x(g')$ for any $g' \neq g$.

The heat bath, or Glauber dynamics, is a Markov chain $(X_n)_{n \in \N}$ where at each step a random element $g \in G$ is chosen. Then the value of $x(g)$ is updated to $-1$ or $1$ with probabilities (proportional to) $e^{-U(y_-)}$ and $e^{-U(y_+)}$, respectively. 
As shown in \cite{ollivier2009ricci, Joulin2010Curvature} for $\beta$ small enough\footnote{A precise statement what `small enough' means can be found in \cite[Example 17]{ollivier2009ricci}.} the chain is contractive with $\pi$ being the invariant distribution.
\end{example}





\begin{example}[Simple Slice Sampling]\label{ex:simple_slice_sampling}
Let $\cX = \R^s$ and let $\rho = d $ be the Euclidean metric.
Assume that the target distribution $\pi$ has a density which satisfies the conditions of \cite[Theorem 2.1]{natarovskii2021quantitative}. 
Roughly speaking, this means that (with additional minor technical requirements) the density of $\pi$ is rotationally invariant. 
We consider a Markov chain $(X_n)_{n \in \N}$ based on (simple) slice sampling, see \cite[Algorithm 1.1]{natarovskii2021quantitative}.
The main idea of the transition mechanism is as follows. 
Given $X_n = x$, one first samples $t$ uniformly from the set $[0, p(x)]$, where $p$ is the (non-normalised) density of $\pi$. Then, $X_{n+1}$ is sampled uniformly from the level set $\{y \in \cX \colon p(y) \geq t\}$.

Indeed, this algorithm targets $\pi$ as invariant distribution and by \cite[Theorem 2.1]{natarovskii2021quantitative} we have that 
\[
\frac{\mathcal{W}(P(x, \cdot ) , P(y, \cdot )) }{d(x,y)} \leq \left( 1- \frac{1}{s+1}\right),
\]
for any $x \neq y \in \R^s$,
where $P$ is the slice sampling Markov kernel and $\mathcal{W}$ is the Wasserstein distance corresponding to $d$. 
Clearly, this implies that $\tau(P) \leq  1- (s+1)^{-1} <1$.
\end{example}

A class of Wasserstein contractive Markov chains are so-called \textit{uniformly ergodic} ones. This is summarised in the upcoming lemma, for a proof see e.g. \cite[Remark 2.1]{rudolf2018pertubation}.
\begin{lem}\label{lem:uniformly_ergodic_implies_contraction}
Let $(X_n)_{n \in \N}$ be a uniformly ergodic Markov chain, that is 
\[
\sup_{x \in \cX}\norm{P^k(x,\cdot) -\pi}_{tv} \leq C \kappa^k,
\]
for some $C\in (0, \infty)$ and $\kappa<1$. Then, $(X_n)_{n \in \N}$ is Wasserstein contractive for $\mathcal{W}$ based on the trivial metric $d (x,y) = 2\cdot\mathds{1}_{\{x \neq y\}}$.
\end{lem}
Uniformly ergodic Markov chains are of particular interest in the Markov chain Monte Carlo literature, as the upcoming three examples demonstrate. 

\begin{example}[Independent Metropolis-Hastings]\label{example:independent_Metropolis_Hastings}
Consider $(X_n)_{n \in \N}$ based on an {independent Metropolis Hastings} algorithm, which is a Markov chain Monte Carlo method to approximate target distributions $\pi$. 
Given $X_n$ a candidate move for $X_{n+1}$ is proposed (independently of $X_n$), and then accepted with a certain probability, see \cite[Algorithm 1]{wang2022exact} for details. 
If the ratio between the density of $\pi$ and the proposal density is uniformly bounded, then the chain is uniformly ergodic \cite{wang2022exact}, whence contractive by Lemma \ref{lem:uniformly_ergodic_implies_contraction}.
\end{example}

\begin{example}[Geodesic Slice Sampling]
We consider a sufficiently regular\footnote{That is, Assumption A and the conditions of Theorem 1 in \cite{hasenpflug2025uniform} should be satisfied.} manifold $\mathscr{M}$. Geodesic slice sampling, see \cite{durmus2025geodesicslicesamplingriemannian}, is a  Markov chain Monte Carlo method to approximately sample from a distribution $\pi$ on $\mathscr{M}$. 
The main idea of the algorithm is similar as in Example \ref{ex:simple_slice_sampling}, however, certain steps are adapted to fit better into the manifold context, we refer to \cite{durmus2025geodesicslicesamplingriemannian}. 
In the setting of \cite[Theorem 1]{hasenpflug2025uniform} the corresponding Markov chain is uniformly ergodic, and therefore contractive by Lemma \ref{lem:uniformly_ergodic_implies_contraction}.
\end{example}

\begin{example}[Stereographic MCMC]
We consider stereographic Markov chain Monte Carlo algorithms, as recently introduced in \cite{yang2024stereographic}.
Let $\cX = \R^s$, and $p \colon \R^s \to \R^+$ be a probability density (w.r.t.~the Lebesgue measure) corresponding to $\pi$.
The main idea of stereographic MCMC is to project $p$ (and hence $\pi$) to the sphere.
In this way it is possible to profit from the geometry of the sphere and obtain chains that mix quickly. 
Since the stereographic projection is bijective, these properties are inherited when projecting back to $\R^s$. 
Indeed, if $p$ is continuous and its tails are not too heavy, then \cite[Theorem 2.1]{yang2024stereographic} shows the uniform ergodicity of their \textit{Stereographic Projection Sampler}. 
By Lemma \ref{lem:uniformly_ergodic_implies_contraction} the corresponding chain is also contractive. 
\end{example}

Another class of Markov chains for which we have a contraction are \textit{$V$-uniformly ergodic} chains. 
These fit well into a `drift and minorisation' framework, see Example \ref{ex:Lyapunov} below, and are also of interest in the MCMC literature. 
For examples of $V$-uniformly ergodic chains, including versions of the Metropolis-Hastings algorithms, see e.g.~\cite[Section 3.4]{roberts2004general}.

The upcoming result, which follows by combining {\cite[Lemma 2.1]{hairer2011yet} and the Kantorovich-Rubinstein duality formula, see e.g. \cite{villani2009optimal}, shows that $V$-uniformly ergodic chains are contractive. 
\begin{lem}\label{lem:V_uniformly_ergodic_implies_contraction}
Let $V \colon \cX \to [1, \infty)$ be $\pi$-integrable, lower semi-continuous and $(X_n)_{n \in \N}$ a $V$-uniformly ergodic Markov chain, that is, for any $k \in \N$,  
\[
\sup_{\vert f \vert \leq V} \left\vert P^kf(x) - \pi(f) \right\vert \leq V(x)C \kappa^k,
\]
for some $C\in (0, \infty)$ and $\kappa<1$ and any $x \in \cX$. Then,  $(X_n)_{n \in \N}$ is Wasserstein contractive for $\mathcal{W}$ based on the metric $d (x,y) = \mathds{1}_{\{x \neq y\}}(V(x) + V(y))$.
\end{lem}

\begin{example}\label{ex:Lyapunov}
	We consider an aperiodic and irreducible Markov chain $(X_n)_{n \in\N}$ with kernel $P$ and invariant distribution $\pi$. Assume there exist a lower semi-continuous and measurable function $V \colon \cX \to [1, \infty)$, a set $S \in \mathcal{F}_\cX$, a probability measure on $\mu$ on $S$, as well as numbers $\lambda, \delta \in (0,1)$ and $b \in (0, \infty)$ such that 
	\begin{itemize}
		\item[1)] for any $x \in \cX $ we have $PV(x) \leq \lambda V(x) + b \mathds{1}_S(x)$ and $\sup_{x \in S} V(x) < \infty$; and
		\item[2)] for any $x \in S$ we have $P(x, A) \geq\delta \mu(A \cap S)$. 
	\end{itemize}
	A function $V$ satisfying 1) is called \textit{Lyapunov function} in the Markov chain literature (see e.g. \cite{meyn2009markov,douc2018markov}). 
	It is known\footnote{We refer e.g. to the appendix of \cite{hofstadler2024convergenceRates} for a proof.} that if $V$ is a Lyapunov function, then so is $V^q$, with $q \in (0,1]$, where in 1) $\lambda $ and $b$ get replaced by  $\lambda^q$ and $b^q$, respectively.  
	
	Let $q < (0,1]$. 
	Define the (lower semi-continuous) metric $d_q(x,y) = \mathds{1}_{\{x \neq y\}} (V^q(x) + V^q(y))$, and let $\mathcal{W}_q$ be the associated Wasserstein distance. 
	Since $V^q$ is also a Lyapunov function and additionally 2) is satisfied, the chain $(X_n)_{n \in \N}$ is $V^q$-uniformly ergodic, c.f.~\cite{Baxendale2005renewal}, and by Lemma \ref{lem:V_uniformly_ergodic_implies_contraction} contractive w.r.t.~$\mathcal{W}_q$.
\end{example}

\begin{example}[MALA]\label{ex:MALA}
Let $\cX = \R^s$, $\rho$ be the Euclidean metric. 
Let $U \colon \cX \to [0, \infty)$ and assume the probability measure $\pi$ has density 
\[
p(x) \propto \exp(-U(x)).
\]
The idea of the Metropolis adjusted Langevin algorithm (MALA) is to make use of the stochastic differential equation (SDE), 
\[
\d Y_t = - \nabla U(Y_t)\d t + \sqrt{2}\d W_t, 
\]
where $W_t$ is a Brownian motion. 
Roughly speaking, one uses an Euler-Maruyama discretisation of this SDE as proposal within a Metropolis-Hastings algorithm, see e.g. \cite{durmus2023MALA} for further details and historical remarks. 

If $U$ is satisfies Assumption 1 and 3 of \cite{durmus2023MALA} and the discretisation step, say $\gamma>0$, is chosen sufficiently small, then \cite[Theorem 1]{durmus2023MALA} implies that the MALA kernel $P_\gamma$ is $V$-uniformly ergodic for 
\[
V(x) = \exp (\eta \norm{x}^2),
\]
with some $\eta >0$, where $\norm{\cdot}$ denotes the Euclidean norm. Hence, the corresponding chain is contractive by Lemma \ref{lem:V_uniformly_ergodic_implies_contraction}.
\end{example}

\begin{example}[No-U-Turn-Sampler]\label{ex:NUTS}
Let $\cX = \R^s$, $\rho$ be the Euclidean metric. 
Let $U \colon \cX \to [0, \infty)$ and assume the probability measure $\pi$ has density 
\[
p(x) \propto \exp(-U(x)).
\]
The No-U-Turn-Sampler (NUTS), see \cite{hoffman2014no}, is a variant of Hamiltonian Monte Carlo, and is currently one of the most popular MCMC algorithms. 
The main idea is to extend $(\cX, \pi)$ 
by adding so-called momentum variables, cf. \cite{hoffman2014no}. 
These  have a physical interpretation and can then be used to efficiently target an extended distribution $\widetilde{\pi}$ having $\pi$ as marginal.  

Recently, it was shown that under suitable assumptions (a variant\footnote{A precise description of the algorithm is beyond the scope of this paper and we refer to \cite{durmus2023convergence} for more details.} of) NUTS is $V$-geometrically ergodic for $V(x) = \exp(\eta \norm{x})$, where $\norm{\cdot} $ is the Euclidean norm and $\eta>0$, cf. \cite[Theorem 17]{durmus2023convergence}.
Thus, the Markov chain corresponding to NUTS is contractive as a consequence of Lemma \ref{lem:V_uniformly_ergodic_implies_contraction}. 
\end{example}

\begin{rem}
Solutions of Poisson's equation for uniformly, or $V$-uniformly ergodic Markov chains were computed in \cite{GLYNN2002hoeffding} and \cite{glynn1996liapounov}, respectively. 
Lemma~\ref{lem:uniformly_ergodic_implies_contraction} and Lemma~\ref{lem:V_uniformly_ergodic_implies_contraction} show that is possible to recover these settings as special cases of Wasserstein contractive chains.
\end{rem}


\begin{rem}
In \cite{eberle2019quantitative} contractions for $\mathcal{W}$ are obtained by transforming the metric $\rho$ into a suitable distance $d$.
This allows to `choose' $d$ based on the Markov chain at hand, cf. Sections 2.1 and 2.2 there, and also Sections 2.4 -- 2.5 for applications to chains induced by Euler schemes, or the Metropolis adjusted Langevin algorithm, respectively. 
\end{rem}



\section{Solving Poission's equation}\label{sec:Solutions}
This section contains our main results about existence and regularity of Poisson's equation for Wasserstein contractive chains. 
First, we consider the case where the forcing function $f$ is Lipschitz w.r.t.~$d$. 
Then, we study the $L^p$ setting, where $p \geq 1$, additionally assuming reversibility.   

Throughout, let $(X_n)_{n \in \N}$ be a Markov chain evolving in $\cX$ with transition kernel $P\colon \cX \times \mathcal{B}(\cX) \to [0,1]$ and unique invariant distribution $\pi$.
\subsection{Solutions for Lipschitz functions}
In this section we solve Poisson's equation \eqref{equ:Poission_equation} for forcing functions $f\colon \cX \to \R$ which are Lipschitz w.r.t.~$d$. 
Recall that $f$ is Lipschitz w.r.t.~$d$ if there exists a constant $L\geq0$ such that 
\begin{equation*}\label{equ:Lipschitz}
\vert f(x) - f(y) \vert \leq L \cdot  d(x,y),
\end{equation*}
for any $x,y \in \cX$. 
Throughout this section, we always refer to the metric $d$ when speaking of Lipschitz functions. 
We denote the space of all measurable Lipschitz functions on $\cX$ by $\mathscr{L}$.
The mapping $\norm{\cdot}_d \colon \mathscr{L} \to [0, \infty)$, given by 
\[
\norm{f}_d = \sup_{x \neq y} \frac{\vert f(x) - f(y)\vert}{d(x,y)},
\] 
is a semi-norm\footnote{That is, $\norm{f}_d \geq 0$,  $\norm{cf}_d = \vert c \vert \cdot \norm{f}_d$ and $\norm{f+g}_d \leq \norm{f}_d + \norm{g}_d$ for any $f,g \in \mathscr{L}$ and $c \in \R$.} on $\mathscr{L}$, cf. \cite{cobzacs2019lipschitz}. 
Note that $\norm{f}_d=0$ if $f$ is constant. 
Hence, $(\mathscr{L}, \norm{\cdot}_d)$ is not a normed space in general. 

The Kantorovich-Rubinstein duality formula, see \cite[Chapter 5]{villani2009optimal}, connects $(\mathscr{L}, \norm{\cdot}_d)$ and $\mathcal{W}$ via 
\[
\mathcal{W}(\mu_1, \mu_2) = \sup_{f \in B_1(\mathscr{L})} \left\vert \int_\cX f(x) \mu_1(\d x) - \int_\cX f(x) \mu_2(\d x) \right\vert,
\]
with $B_1(\mathscr{L}) = \{f \in \mathscr{L} \colon \norm{f}_d \leq 1\}$ being the unit ball in $\mathscr{L}$.


Poisson's equation \eqref{equ:Poission_equation} for forcing function $f$ only makes sense if $\pi(f)$ exists. 
Hence, to study Poisson's equation in the context of Lipschitz functions it is vital that $\pi(f)$ is well-defined for any $f \in \mathscr{L}$.
Following \cite{Joulin2010Curvature} we define the $p$-eccentricity $E_p \colon \cX \to [0,\infty]$, with $p\in [1, \infty)$, by
\[
E_p(x) = \int_\cX d(x,y)^p \pi(\d y).
\]
Note that $E_p(x)$ is well-defined for any $x \in \cX$, but it may be infinite. 
However, if $E(x_0)< \infty$ for some $x_0 \in \cX$, then $E(x)<\infty$ for any $x \in \cX$ by triangle inequality; compare also to \cite{Joulin2010Curvature}.
\begin{lem}\label{lem:integrability_Lipschitz}
Assume that $E_p(x_0)< \infty$ for some $x_0 \in \cX$ and $p \in [1, \infty)$. Then, for any $f \in \mathscr{L}$ we have $\pi(\vert f \vert^p ) < \infty$, and consequently also $\pi(f)$ is well-defined and finite.
\end{lem}

\begin{proof}
For any $f \in \mathscr{L}$ we note that $\vert f \vert^p $ is non-negative and measurable. Observe that $\vert f(x) \vert^p \leq (\vert f(x)- f(x_0) \vert + \vert f(x_0) \vert)^p \leq (d(x,x_0) + \vert f(x_0)\vert)^p$.  
It follows that 
\[
\int_\cX \vert f(x) \vert^p \pi(\d x) \leq  \int_\cX (d(x, x_0) + \vert f(x_0) \vert)^p \pi(\d x)
\leq 2^p\left(  E_p(x_0) + \vert f(x_0) \vert^p\right),
\]
where in the last step we used the inequality $(\alpha+\beta)^p \leq 2^p (\alpha^p+\beta^p)$, which is valid for any $\alpha, \beta \geq 0$.
The statement about $\pi(f)$ is now a simple consequence of Hoelder's inequality (consider $\pi(f \cdot \mathds{1})$).
\end{proof}

Under the conditions of Lemma \ref{lem:integrability_Lipschitz} the right hand side of Poisson's equation \eqref{equ:Poission_equation}, i.e.~$f(x)-\pi(f)$, is well-defined for any $f \in \mathscr{L}$. 
In this case it is no loss of generality to assume $ f \in \mathscr{L}_0 = \{f \in \mathscr{L} \colon \pi(f)=0\}$ which allows to work in a Banach space. 
Although results in this flavour are standard in theory of Lipschitz functions \cite{cobzacs2019lipschitz}, we include a proof for the convenience of the reader. 

\begin{lem}\label{lem:L_0_Banachspace}
Assume that $E_p(x_0) < \infty$ for some $x_0 \in \cX$ and $p \in [1,\infty)$. Then, $(\mathscr{L}_0, \norm{\cdot}_d)$ is a Banach space. 
\end{lem}

\begin{proof}
Let $\xi \in \cX$ be fixed and set $\mathscr{L}_\xi = \{ f \in \mathscr{L} \colon f (\xi)=0\}$.
First we show that $(\mathscr{L}_\xi, \norm{\cdot}_d)$ is a Banach space. 
Namely, let $(f_n)_{n \in \N}$ be a Cauchy sequence in $(\mathscr{L}_\xi, \norm{\cdot}_d)$. Then, there exists a Lipschitz function $f$ with $f(\xi)=0$ such that $\norm{f_n -f}_d \to 0$, see \cite[Theorem 8.1.3]{cobzacs2019lipschitz}.
Additionally, 
\[
\vert f_n(x) - f(x) \vert = \vert f_n(x) - f(x) - (f_n(\xi) -f(\xi)) \vert \leq \norm{f_n - f}_d d(x, \xi),
\]
for any $x \in \cX$. 
Hence, $\lim_{n \to \infty}f_n (x) = f(x)$ for any $x \in \cX$. Thus $f$ is measurable and therefore $f \in \mathscr{L}_\xi$. 

Next we show that $(\mathscr{L}_\xi, \norm{\cdot}_d )$ and $(\mathscr{L}_0, \norm{\cdot}_d)$ are isometrically isomorph. 
To this end, observe that by Lemma~\ref{lem:integrability_Lipschitz} any $f \in \mathscr{L}_\xi$ is $\pi$-integrable.  
We define the mappings $T \colon \mathscr{L}_0 \to \mathscr{L}_\xi$ and $R \colon \mathscr{L}_\xi \to \mathscr{L}_0$ via 
\[
T(f) = f(\cdot) - f(\xi)  \qquad\text{and} \qquad R(h) = h(\cdot) - \pi(h). 
\] 
Note that both, $T$ and $R$, are linear. Furthermore, since $\norm{g}_d$ does not change by adding a constant to $g \in \mathscr{L}$, we deduce 
\[
\norm{Tf }_d = \norm{f}_d \qquad\text{and}\qquad \norm{Rh}_d= \norm{h}_d.
\]
Finally, $T (R(h)) = h$ as well as $R(T(f)) = f$, which shows that the spaces are indeed isometrically isomorph. 
\end{proof}

In what follows we always equip $\mathscr{L}_0$ with the norm $\norm{\cdot}_d$. For the sake of brevity, we simply write $\mathscr{L}_0$ for the (Banach) space $(\mathscr{L}_0, \norm{\cdot}_d)$.
The upcoming lemma is one of our key ingredients to study existence and regularity of solutions to Poisson's equation. 
For a Markov kernel $K\colon \cX \times \mathcal{B}(\cX) \to [0,1]$ it relates $\tau(K)$ to the operator norm of  (the Markov operator) $K$, when acting on $\mathscr{L}_0$.  
In particular, this shows that besides $\tau(K)$ being connected to the ergodic properties of a Markov chain \cite{douc2018markov,rudolf2018pertubation}, or its Ricci curvature \cite{ollivier2009ricci,Joulin2010Curvature,paulin2016mixing}, there is another interpretation\footnote{It should be mentioned that a similar statement as in Lemma \ref{lem:operator_norm_Lipschitz} is also in an early version of \cite{rudolf2018pertubation}. The final (published) version, however, does not use that result.}. 

\begin{lem}\label{lem:operator_norm_Lipschitz}
Assume that $E_p(x_0) < \infty$ for some $x_0 \in \cX$; $p \in [1,\infty)$, and that the Markov kernel $K \colon \cX \times\mathcal{B}(\cX) \to [0,1]$ satisfies $\tau(K)< \infty$.
Then, the mapping $K \colon \mathscr{L}_0 \to \mathscr{L}_0$ given by $f \mapsto Kf$ defines a linear operator and we have 
\[
\norm{K}_{\mathscr{L}_0 \to \mathscr{L}_0} = \tau(K). 
\]  
\end{lem}

\begin{proof}
First, we note that if $\tau(K)<  \infty$, then $K \colon \mathscr{L} \to \mathscr{L}$ indeed defines a linear (and bounded) operator.
Since $\pi$ is the invariant measure, the same is true for $K \colon \mathscr{L}_0 \to \mathscr{L}_0$. 

	
Observe that the semi-norm $\norm{\cdot}_d$ (on $\mathscr{L}$) does not change by adding a constant. That is, for $f \in \mathscr{L}$ and $c \in \R$ we have $\norm{f}_d = \norm{f +c}_d$.
Particularly, by setting $c= \pi(f)$, we deduce from this that 
\[
\sup_{f \in B_1(\mathscr{L}_0)}\norm{K f}_d = \sup_{f \in B_1(\mathscr{L})}\norm{K f}_d,
\]
where $B_1(\mathscr{L})$ and $B_1(\mathscr{L}_0)$ are the unit balls in $\mathscr{L}$ and $\mathscr{L}_0$, respectively. 
Finally, by using the Kantorovich-Rubinstein duality formula\footnote{We refer to \cite[Chapter 5]{villani2009optimal}.}, it is possible to show $ \sup_{f \in B_1(\mathscr{L})}\norm{K f}_d = \tau(K)$, compare to \cite[Chapter 20.3]{douc2018markov}.
\end{proof}




Equipped with Lemma \ref{lem:operator_norm_Lipschitz} we are now ready to look for solutions of Poisson's equation on the Banach space $\mathscr{L}_0$. 




\begin{thm}\label{thm:Solution_Poission_Lipschitz_general}
Assume that $E_p(x_0) < \infty$ for some $x_0 \in \cX$ and $p \in [1,\infty)$ and let $(X_n)_{n \in \N}$ be a Wasserstein contractive Markov chain. 
Then, for any $f \in \mathscr{L}_0$ the space $\mathscr{L}_0$ contains exactly one solution $u = u_f$ of Poisson's equation \eqref{equ:Poission_equation} with forcing function $f$ and 
\begin{itemize}
	\item[i)] $\norm{u}_d \leq \Lambda \norm{f}_d$, as well as
	\item[ii)] $\int_\cX \vert u(x) \vert^p \pi(\d x) \leq 2^p \inf_{y \in \cX}(\vert u(y)\vert^p + \norm{u}_d^p E_p(y)) < \infty$.
\end{itemize}
\end{thm}   

\begin{proof}
Note that all assumptions to apply Lemmas \ref{lem:integrability_Lipschitz}, \ref{lem:L_0_Banachspace} and \ref{lem:operator_norm_Lipschitz} are satisfied, such that 
\[
\norm{P^m}_{\mathscr{L}_0 \to \mathscr{L}_0} = \tau(P^m) <1.
\]
Since $\mathscr{L}_0$ is a Banach space, the Neumann series corresponding to $P$, i.e. $\sum_{k=0}^\infty P^k$, converges. Thus $(\Id-P) \colon \mathscr{L}_0 \to \mathscr{L}_0$ is invertible and  
\[
(\Id - P)^{-1} = \sum_{k=0}^\infty P^k. 
\]
Additionally, $\norm{(\Id - P)^{-1}}_{\mathscr{L}_0 \to \mathscr{L}_0} \leq \Lambda$, with $\Lambda < \infty$ as in Remark \ref{rem:def_Lambda}. 

Given $f \in \mathscr{L}_0$ we observe that for any $u \in \mathscr{L}_0$ we have
\[
(\Id - P)u = f\quad \Longleftrightarrow\quad u = (\Id-P)^{-1}f.
\]
Whence, in the space $\mathscr{L}_0$ the unique solution to Poisson's equation (with forcing function $f$) is given by $u = (\Id-P)^{-1}f$. 

For the rest of the proof fix $f\in \mathscr{L}_0$ and $u= (\Id-P)^{-1}f$.
To see the bound on $\norm{u}_d$ claimed in i), we note that 
\[
\norm{u}_d \leq \norm{(\Id - P)^{-1}}_{\mathscr{L}_0 \to \mathscr{L}_0} \norm{f}_d \leq \Lambda \norm{f}_d.
\]
Observing $E_p(x_0)< \infty$, triangle inequality yields that $E_p(x)< \infty$ for any $x \in \cX$. By the same steps as in the proof of Lemma \ref{lem:integrability_Lipschitz}, we deduce that 
\[
\int_\cX \vert u(x) \vert^p \pi(\d x) \leq 2^p (\vert u(y)\vert^p + \norm{u}_d^p E_p(y))
\]
for arbitrary $y \in \cX$, which implies ii). 
\end{proof}

Theorem \ref{thm:Solution_Poission_Lipschitz_general} establishes the existence of solutions to Poisson's equation. The integrability property ii), however, only provides an upper bound involving function evaluations of $u$.
Under slightly stronger assumptions on the eccentricity it is possible to obtain a refined bound. 
\begin{thm}\label{thm:Solution_Poisson_with_L_p_bound}
Let the setting be as in Theorem \ref{thm:Solution_Poission_Lipschitz_general}, and additionally assume that $\int_\cX \vert E_1(x)\vert^{p_0} \pi(\d x) < \infty$ for some $p_0\in [1, \infty)$. 
Then, for any $f \in \mathscr{L}_0$ the space $\mathscr{L}_0$ contains exactly one solution $u = u_f$ of Poisson's equation with forcing function $f$ which, additionally to the conclusions of Theorem \ref{thm:Solution_Poission_Lipschitz_general}, satisfies
\[
\int_\cX \vert u(x) \vert^{p_0} \pi(\d x) \leq \left( \Lambda \cdot \norm{f}_d \right)^{p_0} \int_\cX \vert E_1(x)\vert^{p_0} \pi(\d x). 
\]
\end{thm}

\begin{proof}
From the proof of Theorem \ref{thm:Solution_Poission_Lipschitz_general} we know that $u $ is given by the Neumann series
\[
u(x) = (\Id-P)^{-1} [f]  (x) = \sum_{k=0}^\infty P^k f(x).
\]
Note that $\vert P^k f(x)\vert \leq \norm{f}_d \mathcal{W}(P^k(x, \cdot), \pi) \leq \norm{f}_d E_1(x) \tau(P^k)$, as a consequence of the Kantorovich-Rubinstein duality formula. Hence 
\begin{equation}\label{equ:pointwise_bound_u}
\vert u(x) \vert^{p_0} \leq \left(\Lambda \cdot \norm{f}_d \right)^{p_0} E_1(x)^{p_0}. 
\end{equation}
Combining \eqref{equ:pointwise_bound_u} with the assumptions on $E_1$ implies the desired result. 
\end{proof}

Using \eqref{equ:pointwise_bound_u} we also obtain the following `$L^\infty$-version' of Theorem \ref{thm:Solution_Poisson_with_L_p_bound}. 

 
\begin{cor}
Let the setting be as in Theorem \ref{thm:Solution_Poission_Lipschitz_general}, and additionally assume that $E_1(x) \leq c < \infty$ for ($\pi$-almost) all $x \in \cX$, with some $c \in (0, \infty)$. Then, for any $f \in \mathscr{L}_0$ the space $\mathscr{L}_0$ contains exactly one solution $u = u_f$ of Poisson's equation with forcing function $f$ which, additionally to the conclusions of Theorem \ref{thm:Solution_Poission_Lipschitz_general}, for ($\pi$-almost) all $x \in \cX$ satisfies
\[
\vert u(x) \vert \leq c \Lambda \norm{f}_d.
\]
\end{cor}



\subsection{Solutions for $L^p$ functions}
In this section we study Poisson's equation in the $L^p$ setting. 
Again, w.l.o.g. we consider centred forcing functions.

Let $p\in [1, \infty)$.
Recall, that the space $L^p(\pi)$ contains all (equivalence classes\footnote{Two functions are considered as equivalent if they are equal $\pi$-almost everywhere.} of) functions $f \colon \cX \to \R$ for which 
\[
\norm{f}_{L^p(\pi)}^p = \int_\cX \vert f(x)\vert^p \pi(\d x) <\infty.
\]
It is well-known that $(L^p(\pi), \norm{\cdot}_{L^{p}(\pi)})$ is a Banach space for any $p\in [1, \infty)$, and that $(L^2(\pi), \norm{\cdot}_{L^2(\pi)})$ is a Hilbert space with inner product
\[
\langle f, g \rangle_{L^2(\pi)} = \int_\cX f(x) g(x) \pi(\d x). 
\]
The linear subspace $L_0^2(\pi) = \{ f \in L^2 \colon \pi(f)=0\}$ of centred $L^2$ functions is again a Hilbert space, also with inner product $\langle \cdot , \cdot  \rangle_{L^2(\pi)}$ and induced norm $\norm{\cdot }_{L^2(\pi)}$. 
Similarly,  for $p \in [1,\infty)$, we define the centred $L^p$ spaces, write $L^p_0(\pi)$, as $L_0^p(\pi) = \{ f \in L^p(\pi) \colon \pi(f) =0\}$.


Any Markov kernel $K$ with invariant distribution $\pi$ defines a linear operator on the $L^p(\pi)$ spaces. 
Indeed, if $f \in L^p(\pi)$, then $Kf \in L^p(\pi)$ and for the corresponding operator norm we know
\[
\norm{K}_{L^p(\pi) \to L^p(\pi)} := \sup_{\norm{f}_{L^p(\pi) \leq 1}} \norm{K f}_{L^p(\pi)} = 1,
\]
where equality is obtained for the constant function $\mathds{1}_\cX$.
We refer e.g. to \cite{rudolf2012explicit} for a proof. 
However, on the $L^p_0$ spaces it is possible to have 
\[
\norm{K}_{L^p_0(\pi) \to L^p_0(\pi)} := \sup_{\norm{f}_{L^p_0(\pi) \leq 1}} \norm{K f}_{L^p(\pi)} < 1.
\]
In particular, if $\norm{K}_{L^2_0(\pi) \to L^2_0(\pi)}<1$, then we say the kernel $K$ (or the corresponding Markov chain) admits a(n) ($L^2$-)\textit{spectral gap}. 
In the MCMC literature chains with a spectral gap are of interest since for those explicit error bounds in the context of numerical integration are known, cf. \cite{rudolf2012explicit, hofstadler2025optimal}. 

The rest of this section is organised as follows. First, we show a general result about solutions of Poisson's equation for chains with a spectral gap, see Section \ref{subsec:general_results_spectral_gap}. Then we provide two results with sufficient conditions for Wasserstein contractive chains to admit a spectral gap, see Section \ref{subsec:sufficient_conditions_spectral_gap}. 

\subsubsection{General chains with a spectral gap}\label{subsec:general_results_spectral_gap}

In this section we consider a generic Markov kernel $K$ on $(\cX, \mathcal{B}(\cX))$ with invariant distribution $\pi$, and assume a spectral gap, i.e. $\norm{K}_{L^2_0(\pi) \to L^2_0(\pi)}< 1$.
We begin with the following auxiliary result, which is a simple combination of \cite[Lemma 3.16]{rudolf2012explicit} and \cite[Proposition 3.17]{rudolf2012explicit}. 

\begin{lem}\label{lem:spectral_gap_implies_L_p_bound}
For the kernel $K$ we have the following bounds
\[
\norm{K^\ell}_{L^p_0(\pi) \to L^p_0(\pi)} \leq 
\begin{cases}
	2^{2/p} \norm{K}_{L^2_0(\pi) \to L^2_0(\pi)}^{2\ell \frac{p-1}{p}} \qquad &\text{ if } p\in (1,2),\\[1.2ex]
	2^{2\frac{p-1}{p}} \norm{K}_{L^2_0(\pi) \to L^2_0(\pi)}^{2\ell/p} &\text{ if } p \in (2,\infty),
\end{cases}
\]
for any  $\ell \in \N$. 
\end{lem} 

\begin{proof}
	For any $p \in (1, \infty)$ define $\Pi \colon L^p(\pi) \to L^p(\pi)$ via $f \mapsto \pi(f)$. 
	Let $\ell \in \N$. 
	Combining \cite[Lemma 3.16]{rudolf2012explicit} and \cite[Proposition 3.17]{rudolf2012explicit} we obtain $	\norm{K^\ell}_{L^p_0(\pi) \to L^p_0(\pi)} \leq \norm{K^\ell-\Pi}_{L^p(\pi) \to L^p(\pi)}$ as well as
	\[
 \norm{K^\ell-\Pi}_{L^p(\pi) \to L^p(\pi)}
	\leq \begin{cases}
		2^{2/p} \norm{K}_{L^2_0(\pi) \to L^2_0(\pi)}^{2\ell \frac{p-1}{p}} \qquad &\text{ if } p\in (1,2),\\
		2^{2\frac{p-1}{p}} \norm{K}_{L^2_0(\pi) \to L^2_0(\pi)}^{2\ell/p} &\text{ if } p \in (2,\infty),
	\end{cases}
	\]
	which yields the desired result. 
\end{proof}

The following theorem is our main result concerning solutions of Poisson's equation in the $L^p_0$ setting for kernels $K$ which admit a spectral gap.
That is, for given $f \in L^p_0(\pi)$ we want to find $u \in L^p_0(\pi)$ such that $Ku - u = f$, understood in the $L^p(\pi)$-sense. 
In \cite[Proposition 21.3.8]{douc2018markov} the existence of $u$ is shown assuming that $\sum_{\ell=0}^\infty \norm{K^\ell f}_{L^p(\pi)}$ converges, which is closely related to a converging Neumann series.
Below we extend the mentioned result by a), showing that a spectral gap is a sufficient condition, and b), providing $L^p(\pi)$ bounds for $u$ in terms of the spectral gap.

\begin{thm}\label{thm:poisson_equation_spectral_gap_general}
Let $p \in (1,\infty)$ and $f \in L^p_0(\pi)$. 
If the Markov kernel $K$ admits a spectral gap, then, the space $L^p_0(\pi)$ contains exactly one solution of Poisson's equation $u \in L^p_0(\pi)$ (with forcing function $f$) and we have 
\[
\norm{u}_{L^p(\pi)}
\leq 
\begin{cases}
	\frac{2^{2/p}}{1- \kappa^{2\frac{p-1}{p}}} \norm{f}_{L^p(\pi)} \qquad&\text{ if } p \in (1,2), \\[1.5ex]
	\frac{1}{1-\kappa} \norm{f}_{L^2(\pi)} &\text{ if } p = 2, \\[1.2ex]
	\frac{2^{2\frac{p-1}{p}}}{1-\kappa^{2/p}}\norm{f}_{L^p(\pi)} &\text{ if } p\in (2, \infty),
\end{cases}
\] 
where $\kappa = \norm{K}_{L^2_0(\pi) \to L^2_0(\pi)}< 1$. 
\end{thm}

\begin{proof}
We split the proof into three cases, where we deal with different values of $p \in(1, \infty)$.

\textbf{Case 1:} Let $p \in (1,2)$. Since $\kappa = \norm{K}_{L^2_0(\pi) \to L^2_0(\pi)} <1$ we have $\kappa^{2\frac{p-1}{p}} <1$. 
By Lemma \ref{lem:spectral_gap_implies_L_p_bound},  
\[
\norm{K^\ell}_{L^p_0(\pi) \to L^p_0(\pi)} \leq 2^{2/p} \left(\kappa^{2\frac{p-1}{p}} \right)^\ell.
\]
Hence the corresponding Neumann series converges and on $L^p_0(\pi)$,  
\[
(\Id - K)^{-1} = \sum_{\ell=0}^\infty K^\ell.
\]
Additionally we have the following bound for the operator norm, 
\[
\norm{(\Id - K)^{-1}}_{L^p_0(\pi) \to L^p_0(\pi)} \leq \frac{2^{2/p}}{1- \kappa^{2\frac{p-1}{p}}}.
\]
Now, $ u = (\Id - K)^{-1} f \Leftrightarrow u - Ku =f$, such that Poisson's equation (with forcing function $f$), admits the unique solution $u$ in $L^p_0(\pi)$, and 
\[
\norm{u}_{L^p(\pi)} \leq \norm{(\Id - K)^{-1}}_{L^p_0(\pi) \to L^p_0(\pi)} \norm{f}_{L^p(\pi)} \leq \frac{2^{2/p}}{1- \kappa^{2\frac{p-1}{p}}} \norm{f}_{L^p(\pi)}.
\] 

\textbf{Case 2:} Let $p=2$. Since $\norm{K}_{L^2_0(\pi) \to L^2_0(\pi)} <1$, we again know that the Neumann series corresponding to $K$ converges and 
\[
\norm{(\Id - K)^{-1}}_{L^2_0(\pi) \to L^2_0(\pi) \to L^2_0(\pi)} \leq \frac{1}{1 - \norm{K}_{L^2_0(\pi) \to L^2_0(\pi)} }. 
\]
Setting $u = (\Id - K)^{-1}f$ we obtain the desired result for $p=2$.

\textbf{Case 3:} Let $p >2$. Here we can use the same arguments as in case 1, but with the corresponding bounds for $p >2$ from Lemma \ref{lem:spectral_gap_implies_L_p_bound}. 
\end{proof}


\subsubsection{Contractive chains with spectral gaps}\label{subsec:sufficient_conditions_spectral_gap}
Here we provide two results which connect contractive chains and spectral gaps. 
These give sufficient conditions which ensure a contractive chain also admits a spectral gap, and whence provide access to Theorem \ref{thm:poisson_equation_spectral_gap_general}. 
However, before we start we require some additional definitions. 

Recall that $(X_n)_{n \in \N}$ is a Markov chain with transition kernel $P$ and invariant distribution $\pi$. 
In the spirit of \cite{ollivier2009ricci} we define
\[
\sigma(x) = \left(\frac{1}{2} \int_\cX \int_\cX d(y,y')^2 P(x, \d y) P(x, \d y') \right)^{1/2},
\]
called the \textit{coarse diffusion} coefficient of $(X_n)_{n \in \N}$, and its $L^2$ norm,
\[
\sigma = \norm{\sigma(\cdot)}_{L^2(\pi)} = \left( \int_\cX \sigma (x)^2 \pi(\d x) \right)^{1/2} . 
\]
Using \cite[Proposition 30]{ollivier2009ricci} we deduce the following result for the case where $d = \rho$, i.e. the metric rendering $\cX$ Polish is the same as in $\mathcal{W}$. 

\begin{lem}[{\cite[Proposition 30]{ollivier2009ricci}}]\label{lem:contraction_implies_spectral_gap}
	Let $d = \rho$.
	Assume $(X_n)_{n \in \N}$ is a Wasserstein contractive Markov chain being reversible w.r.t.~$\pi$.
	If $\sigma < \infty$, then 
	\[
	\norm{P}_{L^2_0(\pi) \to L^2_0(\pi)} \leq \sqrt[m]{\tau(P^m)} <1. 
	\] 
\end{lem}

\begin{proof}
	Proposition 30 of \cite{ollivier2009ricci} implies that the spectral radius of the Markov operator $P^m \colon L^2_0(\pi ) \to L^2_0(\pi)$, say $r(P^m)$, is bounded by $\tau(P^m)<1$. 
	Since $(X_n)_{n \in \N}$ is reversible the operator $P^m$ is self adjoint on $L^2_0(\pi)$ implying
	\[
	\norm{P^m}_{L^2_0(\pi) \to L^2_0(\pi)} = r(P^m) \leq \tau(P^m) <1.
	\]
	Using the formula for the spectral radius of $P$ acting on $L^2_0(\pi)$, say $r(P)$, it follows 
	\[
	\norm{P}_{L^2_0(\pi) \to L^2_0(\pi)} = r(P) = \inf_{\ell \in \N} \norm{P^\ell}^{1/\ell} \leq \tau(P^m)^{1/m}<1.
	\]
\end{proof}

The above lemma provides a criterion for contractive chains to have a spectral gap. 
However, the requirement that $d = \rho$ may not always be true.

The next result, which is \cite[Proposition 2.8]{hairer2014spectral}, see also \cite{rockner2001weak}, shows that $d = \rho$ is not needed if the space of Lipschitz functions (w.r.t.~$d$), say $\mathscr{L}$, is dense in $L^2(\pi)$. 

\begin{lem}[{\cite[Proposition 2.8]{hairer2014spectral}}]\label{lem:contraction_implies_spectral_gap_2}
	Assume $(X_n)_{n \in \N}$ is a Wasserstein contractive Markov chain, which is reversible w.r.t.~$\pi$.
	If $\mathscr{L} \cap L^\infty(\pi)$ is dense in $L^2(\pi)$, then 
	\[
	\norm{P}_{L^2(\pi) \to L^2(\pi)} = \sqrt[m]{\tau(P^m)} < 1. 
	\]
\end{lem}

\begin{proof}
	Define the (reversible/self-adjoint) Markov kernel/operator $K= P^m$, and note that $\tau(K) = \tau(P^m)<1$.
	By \cite[Proposition 2.8]{hairer2014spectral} it follows 
	\[
	\norm{K f}_{L^2(\pi) } \leq \tau(K) \norm{f}_{L^2(\pi)}
	\]
	for any $f \in L^2_0(\pi)$, which clearly implies $\norm{K}_{L^2_0(\pi) \to L^2_0(\pi)} \leq \tau(K)< 1$. 
	By the well-known formula for the spectral radius of $P$, say $\rho$, and since $P$ is self-adjoint on $L^2_0(\pi)$, we arrive at 
	\[
	\norm{P}_{L^2_0(\pi) \to L^2_0(\pi)} = \rho = \inf_{\ell \in \N} \norm{P^\ell}_{L^2_0(\pi) \to L^2_0(\pi)}^{1/\ell} \leq \norm{P^m}_{L^2_0(\pi) \to L^2_0(\pi)}^{1/m} \leq \sqrt[m]{\tau(P^m)}, 
	\]
	which shows the claimed result. 
\end{proof}




\section{Maximal inequalities}\label{sec:applications}
In this section we apply our results to derive maximal inequalities for Wasserstein contractive Markov chains.


Given some integrable $f \colon \cX \to \R$ with $\pi(f)=0$ and $n \in \N$ we define 
\[
S_n^* f := \max_{1 \leq k \leq n} \vert S_k f \vert = \max_{1 \leq k \leq n} \left\vert \sum_{j=1}^k f(X_j) \right\vert.
\]
Estimates for $S_n^*f$ are called \textit{maximal inequalities}. 
For example, Kolmogoroff's inequality bounds $\mathbb{P}[S_n^*f \geq \eps]$ for $f \in L^2(\pi)$ if $(X_n)_{n \in \N}$ is an iid sequence with $X_1 \sim \pi$, cf. \cite[Chapter 9]{dudley2002real}.
Additionally, there are different inequalities due to Doob bounding $\mathbb{P}[S_n^*f \geq \eps]$ or $\mathbb{E}[(S_n^*f)^p]$ provided $(S_n f)_{n \in \N}$ is a sufficiently regular martingale, see e.g. \cite{revuz1999continuous}.  

 
Here we want to derive analogues for Wasserstein contractive Markov chains. 
As before, such a chain is denoted by $(X_n)_{n \in \N}$, its transition kernel by $P$, and its invariant and initial distribution are $\pi$ and $\nu$, respectively.
Recall that the space of measurable Lipschitz functions (w.r.t.~$d$) is $\mathscr{L}$,  and $\mathscr{L}_0\subseteq\mathscr{L}$ contains those $f \in \mathscr{L}$ satisfying $\pi(f)=0$. 
Finally, we remind ourselves about the eccentricities $E_p(x)=\int_\cX d(x,y)^p \pi(\d y)$. 

The main tool in this section is decomposition \eqref{equ:martingale_decomposition}, which shows that 
\[
S_n f = M_n + u(X_1) - u(X_{n+1}),
\]
where $u $ is a solution of Poisson's equation with forcing function $f$. 
Using this identity for any $1 \leq k \leq n$ we obtain 
\begin{equation}\label{equ:maximal_decomposition}
S_n^*f \leq M_n^* + R_n^*,
\end{equation}
with $M_n^* = \max_{1 \leq k \leq n} \vert M_k\vert$ and $R_n^* = \max_{1 \leq k \leq n} \vert u(X_1) - u(X_{k+1}) \vert $.



The remainder of this section is organised as follows. First we prove some preliminary results which provide estimates for $M_n$ under general integrability conditions on $u$.
Then, we prove maximal inequalities for $f \in \mathscr{L}_0$.  

\subsection{Regularity of the martingale $M_n$}
Here we investigate the integrability conditions of $(M_n)_{n \in \N}$.
This is done under general conditions on $u$ and the Markov chain $(X_n)_{n \in \N}$. 
Recall that the underlying probability space is $(\Omega, \mathcal{F}, \mathbb{P}_\nu)$ and that
the filtration $(\mathcal{F}_n)_{n \in \N}$ is defined as $\mathcal{F}_n = \sigma(X_k \colon 1\leq k \leq n+1)$.

Let $\nu \ll \pi$ with Radon-Nikod\'ym derivative $\frac{\d \nu}{\d \pi} \in L^q(\pi)$, for some $q \in (1,\infty]$.
For some generic integrable $f \colon \cX \to \R$ assume $u$ is the solution of Poisson's equation with forcing function $f$. 
Throughout this section $M_n, R_n$ (see \eqref{equ:martingale_decomposition}) correspond to this $f$ and $u$. 

\begin{lem}\label{lem:martingale_bound_1}
Let $s$ be the Hoelder conjugate\footnote{That is, $s \in (1, \infty)$ is the unique number such that $s^{-1} + q^{-1} =1$.} of $q \in (1, \infty)$.
Let $p\geq 1$ and let $r \geq sp$. If $u \in L^{r}(\pi)$, then $(M_n)_{n \in \N}$ is a martingale w.r.t.~$(\mathcal{F}_n)_{n \in \N}$ and for its differences we have
\[
\E_\nu[ \vert M_{n+1} - M_n \vert^p ]^{1/p} \leq 2\norm{\frac{\d \nu}{\d \pi}}_{L^q(\pi)}^{1/p} \norm{u}_{L^r(\pi)}
\]
for any $n \in \N$.
In particular, this implies that $(M_n)_{n \in \N}$ is a $L^p$ martingale. 
\end{lem}

\begin{proof}
We begin by showing the claimed bound on the differences. 
For the sake of brevity write $\Delta_j = u(X_{j+1}) -Pu(X_j) $ such that $M_n = \sum_{j=1}^n \Delta_j$. 

Note that $\mathbb{P}_\nu$ and $\mathbb{P}_\pi$ only differ by the choice of the initial distribution. 
By Hoelder's inequality
\begin{align*}
\E_\nu\left[\vert \Delta_j\vert^p\right]^{1/p} = \E_\pi\left[ \frac{\d \nu}{\d \pi}(X_1) \vert \Delta_j\vert^p\right]^{1/p} 
\leq\norm{\frac{\d \nu}{\d \pi}}_{L^q(\pi)}^{1/p} \E_\pi\left[  \vert \Delta_j\vert^{ps}\right]^{1/(sp)}.
\end{align*}
By assumption $r\geq sp$, such that by definition of $\Delta_j$ we obtain
\begin{align*}
\E_\pi\left[  \vert \Delta_j\vert^{ps}\right]^{1/(sp)}&\leq 
\E_\pi\left[  \vert \Delta_j\vert^{r}\right]^{1/r}\\ 
&\leq 
\E_\pi\left[\vert u(X_{j+1}) \vert^r \right]^{1/r} + \E_\pi\left[\vert Pu(X_{j}) \vert^r \right]^{1/r} 
\leq 2 \norm{u}_{L^r(\pi)},
\end{align*}
where in the last step we used that $\norm{P}_{L^r(\pi) \to L^r(\pi)} =1$. 

It follows that $\E_\nu[\vert M_n \vert] < \infty$ for any $n\in \N$. Moreover, by definition of $M_n$, we also have 
\[
\E [\Delta_n \vert \mathcal{F}_{n-1}] = 0
\]
showing that $(M_n)_{n \in \N}$ is a martingale as claimed. 
\end{proof}

The proof of the previous lemma also transfers to the case $q= \infty$. 
We summarise this in the following lemma, whose proof is left to the reader.
\begin{lem}\label{lem:martingale_bound_2}
Assume $q=\infty$ and let $p \geq 1$. 
If $u \in L^{p}(\pi)$, then $(M_n)_{n \in \N}$ is a martingale w.r.t.~$(\mathcal{F}_n)_{n \in \N}$ and for its differences we have
\[
\E_\nu[ \vert M_{n+1} - M_n \vert^p ]^{1/p} \leq 2\norm{\frac{\d \nu}{\d \pi}}_{L^\infty(\pi)}^{1/p} \norm{u}_{L^p(\pi)}
\]
for any $n \in \N$.
In particular, this implies that $(M_n)_{n \in \N}$ is a $L^p$ martingale. 
\end{lem}





\subsection{Maximal inequalities for Lipschitz functions}
In this section we show bounds for $S_n^* f$, with $f \in \mathscr{L}_0$.
Throughout, we assume that $(X_n)_{n \in \N}$ is contractive, cf. Definition \ref{def:contractive_chain}. 
We define 
\[
\delta := \text{diam}(\cX) = \sup_{x,y \in \cX} d(x,y) . 
\]
We begin by proving results for the case where $\delta<\infty$, and then turn to more general settings. 
\subsubsection{Bounded diameter}

Here we consider the case $\delta<\infty$, implying $E_p(x_0) = \int_{\cX} d(x_0, x)^p \pi(\d x) \leq \delta^p< \infty$ for any $x_0 \in \cX$ and $p\geq 1$. 
Examples include uniform ergodic chains, see Lemma \ref{lem:uniformly_ergodic_implies_contraction}, or compact state spaces $\cX \subseteq \R^s$. 


\begin{prop}\label{prop:Doob_maximal_bounded_diameter_general}
Let $f \in \mathscr{L}_0$. Assume $\frac{\d\nu}{\d \pi}\in L^q(\pi)$ for some $q \in (1, \infty)$, then, for any $p \in [1, \infty)$ and any $t \in (0, \infty)$ we have
\[
\mathbb{P}_\nu\left[  S_n^*f > t \right] \leq \left(\frac{2}{t}\right)^p\left( \E_\nu \left[\vert M_n\vert^p \right] + C^p\norm{f}_d^p \right),
\]
with $C=\delta \cdot \Lambda $ and $M_n$ as in \eqref{equ:martingale_decomposition}.

\end{prop}

\begin{proof}
Observe that $E_r (x) \leq \delta^r$ for any $r \in [1, \infty)$ and $x \in \cX$.


By virtue of Theorem \ref{thm:Solution_Poisson_with_L_p_bound} we find that there exists a solution of Poisson's equation (with forcing function $f$), say $u$, and that $\norm{u}_{L^r(\pi)} \leq C \norm{f}_d$ for any $r \in [1,\infty)$, with $C= \delta \cdot \Lambda$.

Hence Lemma \ref{lem:martingale_bound_1} may be applied with arbitrary $p\in [1, \infty)$ and yields that $(M_n)_{n \in \N}$ is a $L^p$ martingale (on $(\Omega, \mathcal{F}, \mathbb{P}_\nu)$ and w.r.t.~filtration $(\mathcal{F}_n)_{n \in \N}$). For arbitrary $\eps>0$ Doob's maximal inequality, see \cite{revuz1999continuous} - Corollary 1.6 of Chapter II, yields 
\[
\mathbb{P}_\nu\left[ M_n^* > \eps \right] \leq \left(\frac{1}{\eps^p} \E_\nu[\vert M_n\vert^p] \right).
\] 


By the Lipschitz continuity of $u$ it follows 
\[
\vert u(X_k) - u(X_1)\vert \leq \norm{u}_d d(X_k, X_1) \leq C \norm{f}_d .
\]
Hence $R^*_n \leq C \norm{f}_d$ which implies that $\mathbb{P}_\nu[R_n^* > \eps] \leq \left(\frac{C  \norm{f}_d}{\eps}\right)^p$, for any $p \in [1, \infty)$.
Now, the result follows by observing that 
\[
\mathbb{P}_\nu\left[  S_n^*f > t \right] \leq \mathbb{P}_\nu\left[  M_n^*f > t/2 \right] + \mathbb{P}_\nu\left[  R_n^*f > t/2 \right],
\]
which is a consequence of \eqref{equ:maximal_decomposition}.
\end{proof}


The upcoming result can be seen as a version of Doob's maximal inequality, see Corollary 1.6 of Chapter II in \cite{revuz1999continuous}, for contractive chains. 
\begin{thm}\label{thm:Doob_maximal_bounded_diameter}
Assume $\frac{\d\nu}{\d\pi}\in L^q(\pi)$ for some $q \in (1, \infty)$ with Hoelder conjugate $s>1$.
Then, for any $n \in \N$  we have 
\[
\sup_{f \in \mathscr{L}_0; \norm{f}_d \leq 1}
\mathbb{P}_\nu\left[  S_n^*f > t \right] \leq \frac{4}{t^2}\left( 2n\norm{\frac{\d \nu}{\d \pi}}_{L^q(\pi)} C^{2r} + C^2\right),
\]
for any $t>0$,
where $C= \delta\cdot \Lambda$ and $r = 2s$.
\end{thm}

\begin{proof}
	Let $f\in \mathscr{L}_0$ with $\norm{f}_d\leq 1$.  
Note that all assumptions to apply Proposition \ref{prop:Doob_maximal_bounded_diameter_general} with $p=2$ are met.  

Employing Lemma \ref{lem:martingale_bound_1} we know that $(M_n)_{n \in \N}$ is a $L^2$ martingale (on $(\Omega, \mathcal{F}, \mathbb{P}_\nu)$ w.r.t.~filtration $(\mathcal{F}_n)_{n \in \N}$). 
Set $\Delta_j = u(X_{j+1}) -Pu(X_j) $ such that $M_n = \sum_{j=1}^n \Delta_j$.
By the martingale property and Lemma \ref{lem:martingale_bound_1}, 
\[
\E_\nu\left[M_n^2\right]= 
\sum_{j=1}^n\E_\nu \left[\Delta_j^2 \right]
\leq 2n\norm{\frac{\d \nu}{\d \pi}}_{L^q(\pi)} \norm{u}_{L^r(\pi)}^2
\leq 
2n\norm{\frac{\d \nu}{\d \pi}}_{L^q(\pi)} C^{2r},
\]
where im the last step we used Theorem \ref{thm:Solution_Poisson_with_L_p_bound} to bound $\norm{u}_{L^r(\pi)}$.
The desired result is now a consequence of Proposition \ref{prop:Doob_maximal_bounded_diameter_general}.
\end{proof}



Next we present an $L^2$ maximal inequality for contractive chains. 
\begin{thm}\label{thm:L_2_maximal_inequality_bounded_diameter}
Assume $\frac{\d\nu}{\d\pi}\in L^q(\pi)$ for some $q \in (1, \infty)$ with Hoelder conjugate $s>1$.
Then, for any $n \in \N$  we have 
\[
\sup_{f \in \mathscr{L}_0; \norm{f}_d \leq 1} \E_\nu\left[(S_n^*f)^2\right] \leq \left(32C^{2r}n\norm{\frac{\d \nu}{\d \pi}}_{L^q(\pi)}  + 2C^2 \right),
\]
where $C= \delta \cdot \Lambda$ and $r = 2s$.
\end{thm}

\begin{proof}
Let $f\in \mathscr{L}_0$ with $\norm{f}_d \leq 1$.  
As in the proof of Theorem \ref{thm:Doob_maximal_bounded_diameter} we deduce that $(M_n)_{n \in \N}$ is a $L^2$ martingale and 
\[
\E_\nu\left[M_n^2\right] \leq 2n\norm{\frac{\d \nu}{\d \pi}}_{L^q(\pi)} C^{2r},
\]
for any $n \in \N$. 
Additionally, we know that $R_n^* \leq C \norm{f}_d = C $ (see proof of Proposition \ref{prop:Doob_maximal_bounded_diameter_general}).
Using \eqref{equ:maximal_decomposition} and then Doob's $L^2$ inequality, see \cite{revuz1999continuous}, yields  
\begin{align*}
\E_\nu\left[(S_n^*f)^2 \right]^{1/2} &\leq \E_\nu\left[ (M_n^*)^2 \right]^{1/2} + \E_\nu\left[(R_n^*)^2\right]^{1/2} \\
&\leq 4\E_\nu\left[M_n^2\right]^{1/2} + \E_\nu\left[(R_n^*)^2\right]^{1/2}  \leq 4\E_\nu\left[M_n^2\right]^{1/2} + C.
\end{align*}
Combining the inequality $(\alpha+\beta)^2 \leq 2(\alpha^2+\beta^2)$ with the above estimates, the desired result follows. 
\end{proof}

\subsubsection{Distributions with finite moments}
In the previous section we studied the case where $\cX$ has bounded diameter which might not always be the case in applications. Here we prove maximal inequalities under suitable moment assumptions on $E_1(\cdot)$, where one does not necessarily require $\delta < \infty$.
To this end, for $p \in (1, \infty)$, define 
\[
\eps_p  = \norm{E_1(\cdot )}_{L^p(\pi)} = \left( \int_\cX E_1(x)^p \pi(\d x) \right)^{1/p} \in [0, \infty].  
\]
Note that if $\eps_p < \infty $, then also $E_1(x_0) < \infty$ for some $x_0 \in \cX$.
Throughout this section we assume that $(X_n)_{n \in \N}$ is a Wasserstein contractive chain. 

In order prove our main result we shall make use of the following condition about $R_n^*$. 
\begin{ass}\label{ass:bounded_R_n_star}
	
There exists some $\widehat{C} >0$ such that for any forcing function $f \in \mathscr{L}_0$ we have
\[
\E_\nu \left[ \left\vert R_n^*  \right\vert^2  \right] \leq \widehat{C} n \norm{f}_d^2
\] 
for any $n \in \N$, with $R_n^*$ as in \eqref{equ:maximal_decomposition}.  
\end{ass}

\begin{rem}
The uniformity in Assumption \ref{ass:bounded_R_n_star} is restrictive, however, at the end of this section we show that $V$-uniformly ergodic Markov chains (see Lemma \ref{lem:V_uniformly_ergodic_implies_contraction} and the examples below) satisfy such a bound. 
\end{rem}

\begin{thm}\label{thm:Doob_maximal_inequality_finite_moment_assumption}
Let $f \in \mathscr{L}_0$ and $\frac{\d \nu}{\d \pi} \in L^q(\pi)$ for some $q \in (1, \infty)$, with Hoelder conjugate $s \in (1, \infty)$. 
Assume that {$\eps_r<\infty$ for $r=2s$}.
If Assumption~\ref{ass:bounded_R_n_star} is satisfied, then, for any $n \in \N$  we have 
\[
\sup_{f \in \mathscr{L}_0; \norm{f}_d \leq 1}
\mathbb{P}_\nu\left[  S_n^*f > t \right] \leq \frac{4n}{t^2}\left(4 \norm{\frac{\d \nu}{\d \pi}}_{L^q(\pi)} \left(\Lambda \eps_r\right)^2 + \widehat{C} \right),
\]
for any $t>0$.
\end{thm}

\begin{proof}
Let $f \in \mathscr{L}_0$ with $\norm{f}_d\leq 1$ be fixed. 
Due to Theorem \ref{thm:Solution_Poisson_with_L_p_bound} the corresponding solution of Poisson's equation $u$ is in $L^r_0(\pi)$.   
Lemma \ref{lem:martingale_bound_1} implies that the martingale $(M_n)_{n \in \N}$, as specified in \eqref{equ:martingale_decomposition}, has increments, say $(\Delta_j)_{j \in \N}$, with $\Delta_j = u(X_{j+1}) -Pu(X_j) $, which are uniformly bounded in $L^2$, that is 
\[
\E_\nu\left[\Delta_j^2\right] \leq  4\norm{\frac{\d \nu}{\d \pi}}_{L^q(\pi)} \norm{u}_{L^{r}(\pi)}^2.
\]
Doob's maximal inequality (see Corollary 1.6 of Chapter II in \cite{revuz1999continuous}) yields 
\[
\mathbb{P}_\nu\left[ M_n^* > \delta \right] \leq \frac{1}{\delta^2} \E_\nu\left[ M_n^2 \right] = \frac{1}{\delta^2}\E_\nu\left[ \sum_{j=1}^n \Delta_j^2 \right] \leq \frac{4n\norm{\frac{\d \nu}{\d \pi}}_{L^q(\pi)} \norm{u}_{L^{r}(\pi)}^2}{\delta^2},
\]
for any $\delta>0$.

Combining Assumption \ref{ass:bounded_R_n_star} and $\norm{f}_d\leq 1$ we deduce 
\[
\mathbb{P}_\nu\left[ R_n^* > \delta \right] \leq \frac{1}{\delta^2} \E_\nu\left[ (R_n^*)^2 \right] \leq \frac{\widehat{C} \cdot n}{\delta^2}. 
\]
Hence the statement follows by observing that 
\[
\mathbb{P}_\nu\left[ S_n^*f > t \right] \leq \mathbb{P}_\nu\left[ M_n^* > t/2 \right] + \mathbb{P}_\nu\left[ R_n^* > t/2 \right],
\]
and that $\norm{u}_{L^r(\pi)} \leq \Lambda \eps_r $ (see Theorem \ref{thm:Solution_Poisson_with_L_p_bound}). 
\end{proof}

Combining the bounds for $\E_\nu\left[(M_n)^2\right]$ and $\E_\nu\left[ (R_n^*)^2 \right]$ from the proof of Theorem \ref{thm:Doob_maximal_inequality_finite_moment_assumption} with Doob's $L^2$-inequality, cf. \cite{revuz1999continuous}, we obtain the following result. 
\begin{cor}\label{cor:L_2_maximal_inequality_finite_moment_assumption}
Under the same assumptions as in Theorem \ref{thm:Doob_maximal_inequality_finite_moment_assumption}, and with $q,r, \widehat{C}$ as specified there, for any $n \in \N$ we have 
\[
\sup_{f \in \mathscr{L}_0; \norm{f}_d \leq 1}
\E_\nu[(S_n^* f)^2] \leq 
{4n}\left(8 \norm{\frac{\d \nu}{\d \pi}}_{L^q(\pi)} \left(\Lambda\eps_r\right)^2 + \widehat{C} \right).
\]
 
\end{cor}

We conclude this section by showing that Assumption \ref{ass:bounded_R_n_star} is satisfied for the setting considered in Example \ref{ex:Lyapunov}.



\begin{prop}\label{prop:growth_assumption_Lyapunov}
Let $\alpha<1/2$ and $d_\alpha$, $\mathcal{W}_\alpha$, $(X_n)_{n \in \N}$ be as in Example~\ref{ex:Lyapunov}.
Then, Assumption \ref{ass:bounded_R_n_star} is true if $\frac{\d \nu}{\d \pi} \in L^{\beta}(\pi)$ where $\beta$ is the Hoelder conjugate of $1/(2\alpha) >1$.
\end{prop}

\begin{proof}
Let $n \in \N$ be fixed and for $1 \leq j \leq n$ chose $A_j^{(n)} \in \mathcal{F}$ such that on each $A_j^{(n)}$ we have $R_n^* = \vert u(X_1) - u(X_{j+1}) \vert$. 
Note that such a collection of (measurable) sets always exists, and moreover, that it is possible to chose the $A_j^{(n)}$'s disjoint. This implies that $(A_j^{(n)})_{j =1}^n$ is a disjoint partition of $\Omega$. 

It follows that 
\begin{align*}
\E_\nu\left[(R_n^*)^2\right] &= \sum_{j=1}^n \E_\nu\left[ \mathds{1}_{A_j^{(n)}} \vert u(X_1) - u(X_{j+1})\vert^2  \right] \\
&\leq 
\sum_{j=1}^n \E_\nu\left[ \vert u(X_1) - u(X_{j+1})\vert^{2}  \right].
\end{align*}
By Theorem \ref{thm:Solution_Poission_Lipschitz_general} the function $u$ is Lipschitz with constant $C(\tau) \norm{f}_d = {C'}$, where $C(\tau)$ does not depend on $f \in \mathscr{L}_0$. Hence
\begin{align*}
\E_\nu\left[ \vert u(X_1) - u(X_{j+1})\vert^{2}  \right] & \leq (C')^2 \E_\nu[d_\alpha(X_1, X_{j+1})^{2}] \\ 
&= 
(C')^2 \E_\nu[(V^{\alpha}(X_1) + V^{\alpha} (X_{j+1}))^{2}] \\
&\leq 
4(C')^2 \left( \E_\nu[V^{2\alpha}(X_1)] + \E_\nu[ V^{2\alpha} (X_{j+1})]\right) .
\end{align*}
Setting $\theta= \norm{\frac{\d \nu}{\d \pi}}_{L^{\beta}(\pi)}$ we obtain for each $k \in \N$ by Hoelder's inequality,
\begin{align*}
\E_\nu[V^{2\alpha}(X_k)] &\leq \theta \E_\pi[ V (X_{k})]^{2\alpha}.
\end{align*}
Now, $\E_\pi[V(X_k)] = \pi(V) \leq \frac{b}{1-\lambda}$, which follows e.g. by \cite[Proposition 4.24]{hairer2006ergodic}.
Combining all the above equations we arrive at 
\[
\E_\nu\left[(R_n^*)^2\right] \leq \widehat{C} \norm{f}_d^2 n, 
\]
where $\widehat{C} \in (0, \infty)$ is a constant independent of $f$, $u$ and $n$ as required. 
\end{proof}

\begin{rem}
Proposition \ref{prop:growth_assumption_Lyapunov} can be extended to chains for which 
\[
\sup_{j \in \N}
\E_\nu[ d(X_{j+1}, X_1)^2 ] < \infty.
\]
\end{rem}
\subsection{Path-wise stability of No-U-Turn-Samplers}
In this section we demonstrate how maximal inequalities can be used to investigate path-wise convergence of MCMC estimators for $\pi(f)$. 
We focus on the NUTS algorithm, compare Example~\ref{ex:NUTS} and \cite{hoffman2014no,durmus2023convergence}. 
Recently, similar results were obtained in \cite{hofstadler2024convergenceRates} for adaptive MCMC methods, however, the No-U-Turn-Sampler was not studied there.

Recall that in the setting of Example \ref{ex:NUTS} we work in the Euclidean space $\cX = \R^s$. 
Writing $\norm{\cdot}$ for the Euclidean norm, define $V \colon \cX \to [1, \infty)$ by 
\[
V(x) = \exp (\eta \norm{x}),
\]
where $\eta>0$.
We assume that all conditions to apply \cite[Theorem 17]{durmus2023convergence} are satisfied such that NUTS algorithm induces a $V$-uniformly ergodic chain $(X_n)_{n \in \N}$. 
Note that the proof of \cite[Theorem 17]{durmus2023convergence} relies on showing \textit{drift and minorisation conditions}, as in Example~\ref{ex:Lyapunov}. 
Hence, as argued in Example~\ref{ex:Lyapunov}, $(X_n)_{n \in \N}$ is also $V^\alpha$-uniformly ergodic for any $\alpha \in (0,1/2)$. 

For the remainder of this section let $\alpha \in (0, 1/2)$ be fixed. 
By Lemma~\ref{lem:V_uniformly_ergodic_implies_contraction} we obtain that $(X_n)_{n \in \N}$ is contractive w.r.t.~$\mathcal{W}_\alpha$ based on the metric 
\[
d_\alpha(x,y) = \mathds{1}_{\{x \neq y\}} (V^\alpha(x) + V^\alpha(y)). 
\]

The following result characterises the speed of convergence for a single realisation of the NUTS estimator for $\pi(f)$. 
For the sake of exposition we assume that $\frac{\d \nu}{\d \pi}$ exists and is bounded, where $\nu$ is the initial distribution.   
\begin{thm}\label{thm:almost_sure_convergence_NUTS}
Let $f \colon \cX \to \R$ be measurable with $\vert f \vert \leq V^\alpha$. 
Then, there exists $C \colon \Omega \to (0, \infty)$ such that for $\mathbb{P}_\nu$-almost all $\omega \in \Omega$ and any $n \in \N$ we have 
\[
\left\vert \frac{1}{n}\sum_{j=1}^n f(X_j(\omega)) - \pi(f) \right\vert \leq C(\omega) \frac{\log n}{\sqrt{n}} .
\]
\end{thm}

\begin{proof}
Note that $\pi(V) < \infty$ implying that $f \in L^2(\pi)$, hence it is no restriction to assume $\pi(f)=0$. 
Additionally, $\vert f \vert \le V^\alpha$ implies that $f$ is Lipschitz w.r.t.~$d_\alpha$, and Assumption \ref{ass:bounded_R_n_star} is true as a consequence of Proposition~\ref{prop:growth_assumption_Lyapunov}.  



Set $n_k = 2^k$ for and $r_k = \sqrt{k}\log (k)$ for $k \in \N$. Theorem \ref{thm:Doob_maximal_inequality_finite_moment_assumption} implies that 
\[
\mathbb{P}_\nu\left[ S_{n_k}^*f > r_{n_k} t \right] \leq \frac{c}{(tk)^2},
\]
for any $t >0$, with some absolute constant $c>0$.
Employing the lemma of Borel-Cantelli we deduce that $\mathbb{P}_\nu$-almost surely
\[
\lim_{k \to \infty}
\frac{1}{r_{n_k}}S_{n_k}^* f  =0.
\] 
Now, given $n \in\N$, chose $k$ such that $n_{k-1} < n \leq n_k$. 
We observe 
\[
0 \leq \left\vert \frac{1}{r_n} S_n f \right\vert \leq \frac{1}{r_n} S_{n_k}^* f \leq \frac{4}{r_{n_k}} S_{n_k}^*f.
\]
Hence
\[
\lim_{n \to \infty} \left\vert \frac{1}{r_n} S_n f \right\vert =0,
\]
$\mathbb{P}_\nu$-almost surely implying the desired result. 
\end{proof}




\section*{Acknowledgements}
Support by the EPSRC programme grant MaThRad EP/W026899/2 is gratefully acknowledged. 

The author thanks Sam Power and Daniel Rudolf for pointing out useful links to literature, and helpful comments regarding the constants appearing in Section~\ref{sec:Wasserstein_contractive_chains} and the presentation of the paper.


\newcommand{\etalchar}[1]{$^{#1}$}
\begin{thebibliography}{DGK{\etalchar{+}}25}
	
	\bibitem[AB15]{adamczak2015exponential}
	R.~Adamczak and W.~Bednorz.
	\newblock {Exponential concentration inequalities for additive functionals of
		Markov chains}.
	\newblock {\em ESAIM: Probability and Statistics}, 19:440--481, 2015.
	
	\bibitem[Ada08]{adamczak2008tail}
	R.~Adamczak.
	\newblock {A tail inequality for suprema of unbounded empirical processes with
		applications to Markov chains}.
	\newblock {\em Electronic Journal of Probability}, 13:1000 -- 1034, 2008.
	
	\bibitem[AG11]{asmussen2011ergodic}
	S.~Asmussen and P.~W. Glynn.
	\newblock {A new proof of convergence of MCMC via the ergodic theorem}.
	\newblock {\em Statistics \& Probability Letters}, 81(10):1482--1485, 2011.
	
	\bibitem[ALR13]{atchade2013Bayesian}
	Y.~F. Atchadé, N.~Lartillot, and C.~P. Robert.
	\newblock Bayesian computation for statistical models with intractable
	normalizing constants.
	\newblock {\em Brazilian Journal of Probability and Statistics},
	27(4):416--436, 2013.
	
	\bibitem[AM06]{andrieu2006ergodicity}
	C.~Andrieu and {\'E}.~Moulines.
	\newblock On the ergodicity properties of some adaptive {MCMC} algorithms.
	\newblock {\em The Annals of Applied Probability}, 16(3):1462--1505, 2006.
	
	\bibitem[Bax05]{Baxendale2005renewal}
	P.~H. Baxendale.
	\newblock {Renewal theory and computable convergence rates for geometrically
		ergodic Markov chains}.
	\newblock {\em The Annals of Applied Probability}, 15(1B):700 -- 738, 2005.
	
	\bibitem[BMS12]{bolt2012invariance}
	W.~Bo{\l}t, A.~Majewski, and T.~Szarek.
	\newblock {An invariance principle for the law of the iterated logarithm for
		some Markov chains}.
	\newblock {\em Studia Mathematica}, 212:41--53, 2012.
	
	\bibitem[CMN19]{cobzacs2019lipschitz}
	{\c{S}}.~Cobza{\c{s}}, R.~Miculescu, and A.~Nicolae.
	\newblock {\em Lipschitz functions}.
	\newblock Springer, 2019.
	
	\bibitem[DGHR25]{durmus2025geodesicslicesamplingriemannian}
	A.~Durmus, S.~Gruffaz, M.~Hasenpflug, and D.~Rudolf.
	\newblock {Geodesic slice sampling on Riemannian manifolds}.
	\newblock arXiv preprint
	\href{https://arxiv.org/abs/2312.00417}{arXiv23.00417}, 2025.
	
	\bibitem[DGK{\etalchar{+}}25]{durmus2023convergence}
	A.~Durmus, S.~Gruffaz, M.~Kailas, E.~Saksman, and M.~Vihola.
	\newblock {On the convergence of dynamic implementations of Hamiltonian Monte
		Carlo and No U-Turn Samplers}.
	\newblock {\em To appear in Annals of Applied Probabiltiy}, 2025.
	
	\bibitem[DJLV25]{douc2025solvingpoissonequationusing}
	R.~Douc, P.~E. Jacob, A.~Lee, and D.~Vats.
	\newblock {Solving the Poisson equation using coupled Markov chains}.
	\newblock arXiv preprint
	\href{https://arxiv.org/abs/2206.05691}{arXiv:2206.05691}, 2025.
	
	\bibitem[DMPS18]{douc2018markov}
	R.~Douc, {\'E}.~Moulines, P.~Priouret, and P.~Soulier.
	\newblock {\em Markov chains}.
	\newblock Springer, 2018.
	
	\bibitem[Dob56]{dobrushin1956central}
	R.~L. Dobrushin.
	\newblock Central limit theorem for nonstationary Markov chains. I.
	\newblock {\em Theory of Probability \& Its Applications}, 1(1):65--80, 1956.
	
	\bibitem[Dob96]{dobrushin96}
	R.~L. Dobrushin.
	\newblock {\em Lectures on Probability Theory and Statistics: Ecole d'Et{\'e}
		de Probabilit{\'e}s de {S}aint-{F}lour {XXIV}---1994}, chapter Perturbation
	methods of the theory of Gibbsian fields, pages 1--66.
	\newblock Springer Berlin Heidelberg, Berlin, Heidelberg, 1996.
	
	\bibitem[Dud02]{dudley2002real}
	R.~M. Dudley.
	\newblock {\em Real Analysis and Probability}.
	\newblock Cambridge Studies in Advanced Mathematics. Cambridge University
	Press, 2 edition, 2002.
	
	\bibitem[EM19]{eberle2019quantitative}
	A.~Eberle and M.~B. Majka.
	\newblock {Quantitative contraction rates for Markov chains on general state
		spaces}.
	\newblock {\em Electronic Journal of Probability}, 24, 2019.
	
	\bibitem[Ete81]{etemadi1981elementary}
	N.~Etemadi.
	\newblock An elementary proof of the strong law of large numbers.
	\newblock {\em Zeitschrift f{\"u}r Wahrscheinlichkeitstheorie und verwandte
		Gebiete}, 55:119--122, 1981.
	
	\bibitem[GI24]{glynn2024solution}
	P.~W. Glynn and A.~Infanger.
	\newblock {Solution representations for Poisson’s equation, martingale
		structure, and the Markov chain central limit theorem}.
	\newblock {\em Stochastic Systems}, 14(1):47--68, 2024.
	
	\bibitem[GLL25]{glynn2025computableboundssolutionpoissons}
	P.~W. Glynn, N.~Lin, and Y.~Liu.
	\newblock {Computable bounds on the solution to Poisson's equation for general
		Harris chains}.
	\newblock arXiv preprint
	\href{https://arxiv.org/abs/2504.01703}{arXiv:2504.01703}, 2025.
	
	\bibitem[GLLP18]{grama2018limit}
	I.~Grama, R.~Lauvergnat, and {\'E}.~Le~Page.
	\newblock {Limit theorems for Markov walks conditioned to stay positive under a
		spectral gap assumption}.
	\newblock {\em The Annals of Probability}, 46(4):1807--1877, 2018.
	
	\bibitem[GM96]{glynn1996liapounov}
	P.~W. Glynn and S.~P. Meyn.
	\newblock {A Liapounov bound for solutions of the Poisson equation}.
	\newblock {\em The Annals of Probability}, pages 916--931, 1996.
	
	\bibitem[GO02]{GLYNN2002hoeffding}
	P.~W. Glynn and D.~Ormoneit.
	\newblock Hoeffding's inequality for uniformly ergodic {Markov} chains.
	\newblock {\em Statistics \& Probability Letters}, 56(2):143--146, 2002.
	
	\bibitem[Hai18]{hairer2006ergodic}
	M.~Hairer.
	\newblock {Ergodic properties of Markov processes}.
	\newblock {\em \href{https://www.hairer.org/notes/Markov.pdf}{Lecture notes}},
	2018.
	
	\bibitem[Has25]{hasenpflug2025uniform}
	M.~Hasenpflug.
	\newblock Uniform ergodicity of geodesic slice sampling.
	\newblock arXiv preprint
	\href{https://arxiv.org/abs/2510.06748}{arXiv:2510.06748}, 2025.
	
	\bibitem[HG14]{hoffman2014no}
	M.~D. Hoffman and A.~Gelman.
	\newblock {The No-U-Turn sampler: adaptively setting path lengths in
		Hamiltonian Monte Carlo}.
	\newblock {\em J. Mach. Learn. Res.}, 15(1):1593--1623, 2014.
	
	\bibitem[HL25]{herve2025computable}
	L.~Herv{\'e} and J.~Ledoux.
	\newblock Computable bounds for solutions to poisson’s equation and
	perturbation of markov kernels.
	\newblock {\em Bernoulli}, 31(3):2302--2322, 2025.
	
	\bibitem[H{\L}RR26]{hofstadler2024convergenceRates}
	J.~Hofstadler, K.~{\L}atuszy{\'n}ski, G.~Roberts, and D.~Rudolf.
	\newblock {Almost sure convergence rates of adaptive increasingly rare Markov
		chain Monte Carlo}.
	\newblock {\em Stochastic Processes and their Applications}, 196:104905, 2026.
	
	\bibitem[HM11]{hairer2011yet}
	M.~Hairer and J.~C. Mattingly.
	\newblock Yet another look at {H}arris' ergodic theorem for {M}arkov {c}hains.
	\newblock In R.~Dalang, M.~Dozzi, and F.~Russo, editors, {\em Seminar on
		Stochastic Analysis, Random Fields and Applications VI}, pages 109--117,
	Basel, 2011. Springer Basel.
	
	\bibitem[Hof25]{hofstadler2025optimal}
	J.~Hofstadler.
	\newblock {Optimal convergence rates of MCMC integration for functions with
		unbounded second moment}.
	\newblock {\em Journal of Applied Probability}, 62(3):1069–1075, 2025.
	
	\bibitem[HSV14]{hairer2014spectral}
	M.~Hairer, A.~M. Stuart, and S.~J. Vollmer.
	\newblock Spectral gaps for a {M}etropolis–{H}astings algorithm in infinite
	dimensions.
	\newblock {\em The Annals of Applied Probability}, 24(6):2455--2490, 2014.
	
	\bibitem[JO10]{Joulin2010Curvature}
	A.~Joulin and Y.~Ollivier.
	\newblock {Curvature, concentration and error estimates for Markov chain Monte
		Carlo}.
	\newblock {\em The Annals of Probability}, 38(6):2418 -- 2442, 2010.
	
	\bibitem[Kol33]{kolmogoroff1933grundbegriffe}
	A.~Kolmogoroff.
	\newblock {\em Grundbegriffe der Wahrscheinlichkeitsrechnung}.
	\newblock Springer, 1933.
	
	\bibitem[LV24]{laitinen2024invitation}
	P.~Laitinen and M.~Vihola.
	\newblock {An invitation to adaptive Markov chain Monte Carlo convergence
		theory}.
	\newblock arXiv preprint
	\href{https://arxiv.org/abs/2408.14903}{arXiv:2408.14903}, 2024.
	
	\bibitem[Mai78]{maigret1978theoreme}
	N.~Maigret.
	\newblock {Th{\'e}or{\`e}me de limite centrale fonctionnel pour une cha{\^\i}ne
		de Markov r{\'e}currente au sens de Harris et positive}.
	\newblock In {\em Annales de l'institut Henri Poincar{\'e}. Section B. Calcul
		des probabilit{\'e}s et statistiques}, volume~14, pages 425--440, 1978.
	
	\bibitem[Mia14]{MIASOJEDOW2014hoeffding}
	B.~Miasojedow.
	\newblock {Hoeffding’s inequalities for geometrically ergodic Markov chains
		on general state space}.
	\newblock {\em Statistics \& Probability Letters}, 87:115--120, 2014.
	
	\bibitem[MT09]{meyn2009markov}
	S.~P. Meyn and R.~L. Tweedie.
	\newblock {\em Markov chains and stochastic stability}.
	\newblock Springer Science \& Business Media, 2009.
	
	\bibitem[Nev72]{neveu1972potentiel}
	J.~Neveu.
	\newblock {Potentiel markovien r{\'e}current des cha{\^\i}nes de Harris}.
	\newblock In {\em Annales de l'institut Fourier}, volume~22, pages 85--130,
	1972.
	
	\bibitem[NRS21]{natarovskii2021quantitative}
	V.~Natarovskii, D.~Rudolf, and B.~Sprungk.
	\newblock {Quantitative spectral gap estimate and Wasserstein contraction of
		simple slice sampling}.
	\newblock {\em The Annals of Applied Probability}, 31(2):806 -- 825, 2021.
	
	\bibitem[Num91]{nummelin1991poisson}
	E~Nummelin.
	\newblock {On the Poisson equation in the potential theory of a single kernel}.
	\newblock {\em Mathematica Scandinavica}, pages 59--82, 1991.
	
	\bibitem[ODM24]{durmus2023MALA}
	A.~Oliviero-Durmus and É. Moulines.
	\newblock {On geometric convergence for the Metropolis-adjusted Langevin
		algorithm under simple conditions}.
	\newblock {\em Biometrika}, 111(1):273--289, 2024.
	
	\bibitem[Oll09]{ollivier2009ricci}
	Y.~Ollivier.
	\newblock {Ricci curvature of Markov chains on metric spaces}.
	\newblock {\em Journal of Functional Analysis}, 256(3):810--864, 2009.
	
	\bibitem[Pau16]{paulin2016mixing}
	D.~Paulin.
	\newblock {Mixing and concentration by Ricci curvature}.
	\newblock {\em Journal of Functional Analysis}, 270(5):1623--1662, 2016.
	
	\bibitem[Pol84]{pollard1984stochastic}
	D.~Pollard.
	\newblock {\em Convergence of Stochastic Processes}.
	\newblock Springer New York, NY, 1984.
	
	\bibitem[QBG25]{qu2025deeplearningpoisson}
	Y.~Qu, J.~Blanchet, and P.~W. Glynn.
	\newblock {Deep Learning for Markov chains: Lyapunov functions, Poisson’s
		equation, and stationary}.
	\newblock arXiv preprint
	\href{https://arxiv.org/abs/2508.2508.17737}{arXiv:2508.2508.17737}, 2025.
	
	\bibitem[Rev84]{revuz1984markov}
	D.~Revuz.
	\newblock {\em Markov Chains}, volume~11.
	\newblock North-Holland Publishing, second edition, 1984.
	
	\bibitem[RR04]{roberts2004general}
	G.~O. Roberts and J.~S. Rosenthal.
	\newblock {General state space Markov chains and MCMC algorithms}.
	\newblock {\em Probability Surveys}, 1(none):20 -- 71, 2004.
	
	\bibitem[RS18]{rudolf2018pertubation}
	D.~Rudolf and N.~Schweizer.
	\newblock {Perturbation theory for Markov chains via Wasserstein distance}.
	\newblock {\em Bernoulli}, 24(4A):2610 -- 2639, 2018.
	
	\bibitem[Rud12]{rudolf2012explicit}
	D.~Rudolf.
	\newblock {Explicit error bounds for Markov chain Monte Carlo}.
	\newblock {\em Dissertationes Math.}, 485, 2012.
	
	\bibitem[RW01]{rockner2001weak}
	M.~R{\"o}ckner and F.-Y. Wang.
	\newblock {Weak Poincar{\'e} inequalities and L2-convergence rates of Markov
		semigroups}.
	\newblock {\em Journal of Functional Analysis}, 185(2):564--603, 2001.
	
	\bibitem[RY99]{revuz1999continuous}
	D.~Revuz and M.~Yor.
	\newblock {\em Continuous martingales and Brownian motion}, volume 293.
	\newblock Springer Science \& Business Media, 3rd edition, 1999.
	
	\bibitem[SV10]{saksman2010ergodicity}
	E.~Saksman and M.~Vihola.
	\newblock {On the ergodicity of the adaptive Metropolis algorithm on unbounded
		domains}.
	\newblock {\em Annals of Applied Probability}, 20(6):2178--2203, 2010.
	
	\bibitem[Vil09]{villani2009optimal}
	C.~Villani.
	\newblock {\em Optimal transport: old and new}, volume 338 of {\em Grundlehren
		der mathematischen Wissenschaften}.
	\newblock Springer, 2009.
	
	\bibitem[Wan22]{wang2022exact}
	G.~Wang.
	\newblock Exact convergence analysis of the independent {Metropolis-Hastings}
	algorithms.
	\newblock {\em Bernoulli}, 28(3):2012--2033, 2022.
	
	\bibitem[Y{\L}R24]{yang2024stereographic}
	J.~Yang, K.~{\L}atuszy{\'n}ski, and G.~O. Roberts.
	\newblock {Stereographic Markov chain Monte Carlo}.
	\newblock {\em The Annals of Statistics}, 52(6):2692--2713, 2024.
	
\end{thebibliography}


\end{document}
